\chapter{Recursors with a built-in seed}
\label{prov:ch:seed}
\label{seed:ch}

The iterator $\It{n}{a}{u}{h}$ of \cref{base:sec:terms} takes its initial value
$a$ as an independent argument and returns it unchanged at the counter $0$.
This chapter treats three arithmetic bases in which the value at $0$ is tied to
other data of the recursor.
\begin{itemize}[nosep]
\item The \emph{successor-seeded} recursor $R(n,a;u.h)$ starts at the successor
of its initializer, $R(0,a;u.h)=\Suc a$. The successor is then the abbreviation
$\Suc a\defeq R(0,a;u.0)$, and the term language is $x\mid 0\mid R(n,a;u.h)$.
\item The \emph{tied-seed} recursor $B(n;u,v.h)$ has no initializer. Its update
sees the old counter $u$ and the old state $v$, and its value at $0$ is
$\Suc(h\subst{0}{u}\subst{0}{v})$.
\item The \emph{self-seeded} recursor $F(n;v.h)$ has no initializer and its
update sees only the old state; its value at $0$ is $\Suc(h\subst{0}{v})$.
\end{itemize}
The hypotheses differ from base to base, and \cref{seed:tab:hyp} records which
properties are axioms and which are derived. In every case the logical base is
either ordinary intuitionistic first-order logic with equality or the native
binders $\Q{y}$, $\D{x}{y}$ of \cref{base:sec:kernel}.

\begin{table}[ht]
\centering\small
\begin{tabular}{@{}llp{0.27\linewidth}p{0.31\linewidth}@{}}
\toprule
Theory & Logic & Axioms beyond the two recursor equations & Derived\\
\midrule
$\mathsf{N}_R$ & native & none (the eight rule forms of $\Nat$) &
predecessor, iterator, hence everything derived in $\Nat$\\
$\mathsf{T}_R$ & ordinary & coverage, $0=\Suc 0\to\bot$, induction for literal
equations & predecessor, injectivity of $\Suc$, iterator\\
$\mathsf{T}_{B\varepsilon}$ & ordinary & Boolean equality test $\varepsilon$
(range, reflection), $0=\Suc 0\to\bot$, induction for literal equations &
predecessor, injectivity, quantifier-free induction, coverage\\
$\mathsf{T}_{Bd}$ & ordinary & two equations for $d$, coverage,
$0=\Suc 0\to\bot$, induction for literal equations & predecessor, injectivity,
reflection of equality, quantifier-free induction\\
$\mathsf{N}_{Bd}$ & native & two equations for $d$ & the same, and the
ordinary connectives\\
$\mathsf{T}_F$ & ordinary & injectivity of $\Suc$, $0=\Suc 0\to\bot$, Boolean
equality test, induction for literal equations & coverage, predecessor (as a
finite-case graph)\\
$\mathsf{T}_F^{-}$ & ordinary & as $\mathsf{T}_F$ without injectivity &
coverage, a right predecessor, $\Suc^k a\neq a$ for each fixed $k\geq1$\\
\bottomrule
\end{tabular}
\caption{Axioms and derived properties of the theories of this chapter.}
\label{seed:tab:hyp}
\end{table}

All recursors are staged: the body $h$ of a recursor term is a term built
earlier, and each fixed body may be read as a function symbol whose arguments
are the counter, the initializer where there is one, and the free variables of
the body. No theory of this chapter quantifies over bodies or has an evaluator
for them. Substitution is capture-avoiding. The numerals are $1\defeq\Suc 0$,
$2\defeq\Suc 1$, and so on; $\Suc^k t$ denotes $k$ written applications of
$\Suc$, and $\bar k\defeq\Suc^k 0$.

%======================================================================
\section{The successor-seeded recursor}
\label{seed:sec:succ}

This section draws on \pkg{scalar\_successor}, \pkg{transient\_seed} and
\pkg{successor\_seed\_completion}, three write-ups of the same recursor.

\subsection{Terms, equations and the two theories}

\begin{definition}\label{seed:def:succ}
The $R$-terms are
\[
 t \mathrel{::=} x \mid 0 \mid R(n,a;u.h),
\]
where $R(n,a;u.h)$ binds $u$ in $h$ and nowhere else. The successor is the
abbreviation $\Suc a\defeq R(0,a;u.0)$. The two computation equations are
\begin{align}
 R(0,a;u.h)&=\Suc a, \label{seed:eq:R0}\\
 R(\Suc n,a;u.h)&=h\subst{R(n,a;u.h)}{u}. \label{seed:eq:RS}
\end{align}
With the abbreviation expanded, \eqref{seed:eq:R0} reads
$R(0,a;u.h)=R(0,a;u.0)$ and the counter on the left of \eqref{seed:eq:RS} is
$R(0,n;u.0)$; no successor constructor is hidden.
\begin{enumerate}[label=(\alph*),nosep]
\item $\mathsf{N}_R$ is the native calculus with the eight rule forms of
$\Nat$ (\cref{base:sec:kernel}) over the $R$-terms: \rulename{Hyp},
\eqref{seed:eq:R0} and \eqref{seed:eq:RS} as the two computation equations,
\rulename{Tr}, \rulename{QE}, \rulename{DE}, \rulename{DI} by cases and
\rulename{QI} by progress for literal equations. The contradiction is
$0=\Suc 0$.
% ed.: the three packages state the calculus with formula transport and with
% progress on the paired binder; the remark on rule forms at the end of the
% section shows that the same sequents are derivable.
\item $\mathsf{T}_R$ is ordinary intuitionistic first-order logic with equality
over the $R$-terms, with \eqref{seed:eq:R0}, \eqref{seed:eq:RS}, coverage
$\forall x\,(x=0\lor\exists y\,x=\Suc y)$, separation $0=\Suc 0\to\bot$, and
induction for literal equations $E$ with side assumptions:
from $\Gamma\proves E(0)$ and $\Gamma,E(z)\proves E(\Suc z)$, with $z$ fresh,
infer $\Gamma\proves E(t)$.
\end{enumerate}
\end{definition}

Neither theory has an axiom of successor injectivity, a predecessor, a
selector or an initializer. The term language is that of the seeded recursor
of \cref{ctr:def:languages}, where it is studied on canonical terms with
normalizing substitution; here the terms are raw and both computation
equations are rules. Weakening, substitution and cut
(\cref{base:lem:struct,base:lem:cut}) hold for $\mathsf{N}_R$ with the same
proofs: the only cases that mention a term constructor are those of the two
computation equations, and \eqref{seed:eq:R0}, \eqref{seed:eq:RS} are closed
under substitution like \eqref{base:eq:I0}, \eqref{base:eq:IS}. In the same
way every result of \cref{base:sec:structural} whose proof does not use a
computation equation holds for $\mathsf{N}_R$, and for the other native
calculi of this chapter; these results are cited below for such calculi
without further comment.

Throughout the section
\[
 \sel^{+}(n,a,b)\defeq R(n,a;u.b)\quad(u\notin\FV(b)),\qquad
 Z(n)\defeq\sel^{+}(n,0,0).
\]
The computation equations give
\begin{equation}\label{seed:eq:selplus}
 \sel^{+}(0,a,b)=\Suc a,\qquad \sel^{+}(\Suc n,a,b)=b,\qquad
 Z(0)=1,\qquad Z(\Suc n)=0 .
\end{equation}
The zero branch of $\sel^{+}$ adds one to the supplied value, so $\sel^{+}$
is not the selector $\sel$ of \cref{base:sec:terms}.

\subsection{Equality, cases and induction before the predecessor}

\begin{lemma}\label{seed:lem:succ-basic}
$\mathsf{N}_R$ derives $t=t$ for every term $t$, symmetry and transitivity of
equality, congruence $s=t\proves h(s)=h(t)$ for every term template $h$, and
transport of literal equations: from $s=t$ and $l(s)=r(s)$ infer
$l(t)=r(t)$. It admits the rules
\[
 \infrule{\Gamma\proves B(0)\qquad\Gamma\proves B(\Suc z)}{\Gamma\proves B(t)}
 \ (\textsc{nc}),\qquad
 \infrule{\Gamma\proves E(0)\qquad\Gamma,E(z)\proves E(\Suc z)}
         {\Gamma\proves E(t)}\ (\textsc{ei}),
\]
where $B$ is any native formula, $E$ is a literal equation and $z$ is fresh
for $\Gamma$, $t$ and the parameters of the matrix. Under the hypothesis
$0=\Suc 0$ it derives every equation.
\end{lemma}

\begin{proof}
Let $r\defeq R(\Suc 0,0;u.t)$ with $u\notin\FV(t)$. Equation
\eqref{seed:eq:RS} is $r=t$. Rule \rulename{Tr} (from $s=t$ and $h(s)=u$ infer
$u=h(t)$) with the source equation $r=t$, the template $h(v)\defeq v$ and the
premise $r=t$ gives $t=t$. This replaces the two instances of
\eqref{base:eq:I0} in \cref{base:lem:eq}(i). The proofs of
\cref{base:lem:eq}(ii)--(v) use \rulename{Tr} and reflexivity only, and give
symmetry, transitivity, congruence and transport of equations. Rule
\textsc{nc} is \cref{base:cor:cases} and rule \textsc{ei} is
\cref{base:lem:eqind}; their proofs use \rulename{DI}, \rulename{DE},
\rulename{QI}, \rulename{QE} and \rulename{Hyp} and no computation equation.

Under $0=\Suc 0$, congruence gives
$\sel^{+}(0,0,s)=\sel^{+}(\Suc 0,0,s)$, that is $\Suc 0=s$ by
\eqref{seed:eq:selplus}, for every $s$. Hence $s=t$ for all $s$, $t$.
\end{proof}

In $\mathsf{T}_R$ reflexivity and replacement are rules of the logic,
\textsc{nc} follows from coverage, \textsc{ei} is the induction rule, and
every formula follows from $0=\Suc 0$ by separation.

\begin{lemma}[Uniqueness of iteration]\label{seed:lem:succ-unique}
Let $u$ be a variable not free in $\Gamma$. If
$\Gamma\proves h=k$, then $\Gamma\proves R(n,a;u.h)=R(n,a;u.k)$ for all terms
$n$, $a$. More generally, if $\Gamma\proves\Suc a=\Suc b$ and
$\Gamma\proves h=k$, then $\Gamma\proves R(n,a;u.h)=R(n,b;u.k)$.
\end{lemma}

\begin{proof}
By \textsc{ei} on a fresh counter $z$ for the literal equation
$R(z,a;u.h)=R(z,b;u.k)$. At $0$ both sides are computed by
\eqref{seed:eq:R0} to $\Suc a$ and $\Suc b$. At $\Suc z$ the left side is
$h\subst{R(z,a;u.h)}{u}$, which equals $h\subst{R(z,b;u.k)}{u}$ by the
induction hypothesis and congruence, and this equals
$k\subst{R(z,b;u.k)}{u}$ by the instance of $\Gamma\proves h=k$ at
$u\defeq R(z,b;u.k)$, which exists because $u$ is not free in $\Gamma$.
Then substitute $n$ for $z$.
\end{proof}

Two bodies with different syntax are thus identified by a proof, not by a
primitive extensionality rule. The lemma is the counterpart for $R$ of
\cref{base:lem:uniq}, proved by the same equation induction with
\eqref{seed:eq:R0}, \eqref{seed:eq:RS} in place of the iterator equations.

\begin{lemma}[Replacement]\label{seed:lem:succ-leibniz}
Write $A\approx A'$ if $A$ and $A'$ are interderivable over the empty
context; by weakening and substitution they are then interderivable over
every context and under every substitution.
\begin{enumerate}[label=(\roman*),nosep]
\item If $\proves t=t'$ is derivable in the empty context, and $A'$ arises
from $A$ by replacing occurrences of $t$ by $t'$, also in the scope of binders
that bind variables of $t$, then $A\approx A'$, in $\mathsf{N}_R$ and in
$\Nat$.
\item The Leibniz rule (\cref{base:thm:leibniz}) and transport in a formula
template (\cref{base:cor:tr}) hold in $\mathsf{N}_R$.
\end{enumerate}
\end{lemma}

\begin{proof}
(ii) The proofs of \cref{base:lem:cong,base:lem:bindercong,base:thm:leibniz}
use the equality rules, \rulename{QE}, \rulename{DE}, \rulename{DI} and the
structural properties, and no computation equation.

(i) By induction on $A$. At a literal equation the bound variables of $A$
that occur in $t$ are free, and transport of equations along $t=t'$ applies.
At a binder the matrices are interderivable over the empty context by the
induction hypothesis, and \cref{base:lem:bindercong} applies, since no
variable is free in the empty context. Statement (i) is not an instance of
\cref{base:thm:leibniz}, whose substitution is capture-avoiding.
\end{proof}

\subsection{Boolean terms and contextual case analysis}

The proof of the predecessor theorem splits on the value of a two-valued term
inside an induction step, where the induction hypothesis is in the context and
mentions that term. Rule \textsc{nc} does not allow this: its premises carry no
equations $t=0$, $t=\Suc z$, and its context may not mention the case
variable. In $\mathsf{T}_R$ the split is a disjunction elimination. In
$\mathsf{N}_R$ it is derived as follows, before any selector with unshifted
branches exists.

\begin{lemma}[Functional range]\label{seed:lem:succ-range}
For every formula template $K(w)$ and every term $r$: if
$\Gamma\proves K(1)$ and $\Gamma\proves K(0)$ then
$\Gamma\proves K(Z(r))$. The context may mention $r$.
\end{lemma}

\begin{proof}
Apply \textsc{nc} to the template $w\mapsto K(Z(w))$ with $w$ fresh. Its zero
premise $K(Z(0))$ follows from $K(1)$ and $Z(0)=1$, and its successor premise
$K(Z(\Suc z))$ from $K(0)$ and $Z(\Suc z)=0$, by
\cref{seed:lem:succ-leibniz}(ii). Then instantiate $w$ by $r$.
\end{proof}

Define the parity terms
\[
 U(n)\defeq R(n,0;u.Z(u)),\qquad N(n)\defeq Z(U(n)).
\]

\begin{lemma}[Parity and the Boolean decoder]\label{seed:lem:succ-decoder}
$\mathsf{N}_R$ derives
\begin{gather*}
 Z(Z(Z(r)))=Z(r),\qquad N(0)=0,\qquad N(\Suc n)=Z(N(n)),\qquad
 N(\Suc\Suc n)=N(n),\\
 N(Z(r))=Z(r),\qquad N(\Suc\Suc Z(r))=Z(r).
\end{gather*}
\end{lemma}

\begin{proof}
The first equation is \cref{seed:lem:succ-range} for the template
$Z(Z(w))=w$: its instances are $Z(Z(1))=Z(0)=1$ and $Z(Z(0))=Z(1)=0$.
By \eqref{seed:eq:R0}, $U(0)=1$, so $N(0)=Z(1)=0$. By \eqref{seed:eq:RS},
$U(\Suc n)=Z(U(n))$, so $N(\Suc n)=Z(Z(U(n)))=Z(N(n))$, and
$N(\Suc\Suc n)=Z(Z(Z(U(n))))=Z(U(n))=N(n)$. In particular $N(1)=Z(0)=1$.
The equation $N(Z(r))=Z(r)$ is \cref{seed:lem:succ-range} for the template
$N(w)=w$, with instances $N(1)=1$ and $N(0)=0$. The last equation follows
from the fourth and the fifth.
\end{proof}

\begin{definition}\label{seed:def:boolean}
A term $b$ is \emph{Boolean in $\Gamma$} if $\Gamma\proves b=Z(r)$ for some
term $r$. The numerals $0$ and $1$ are Boolean in every context, since
$Z(1)=0$ and $Z(0)=1$; so is every term $N(n)$.
\end{definition}
% ed.: transient_seed defines a Boolean term syntactically, as a term Z(r)
% or a numeral 0, 1 "identified with" Z(1), Z(0); provable equality with
% some Z(r) is what its proofs use and makes the identification explicit.

If $b$ and $c$ are Boolean in $\Gamma$, then
$\Gamma,\Suc b=\Suc c\proves b=c$: apply the template $w\mapsto N(\Suc w)$ to
the hypothesis and use the last equation of \cref{seed:lem:succ-decoder} on
both sides. This cancels a successor on Boolean terms only; it is not
injectivity of $\Suc$.

\begin{lemma}[Restricted implication]\label{seed:lem:succ-imp}
Let $b$, $c$ be Boolean in $\Gamma$, let $s=t$ be an equation, and put
\[
 \mathcal I(b=c;\,s=t)\defeq
 \D{x}{y}\bigl(\sel^{+}(y,b,s)=\sel^{+}(y,c,t)\bigr)
\]
with $x$, $y$ fresh. Then
\begin{enumerate}[label=(\roman*),nosep]
\item if $\Gamma,b=c\proves s=t$, then $\Gamma\proves\mathcal I(b=c;\,s=t)$;
\item if $\Gamma\proves\mathcal I(b=c;\,s=t)$ and $\Gamma\proves b=c$, then
$\Gamma\proves s=t$.
\end{enumerate}
\end{lemma}

\begin{proof}
(i) Use \rulename{DI} by cases with witness $0$. The local assumption is the
matrix at $y=0$, which by \eqref{seed:eq:selplus} is $\Suc b=\Suc c$; the
decoder gives $b=c$, and cutting in the given derivation gives $s=t$. By
\eqref{seed:eq:selplus} this is the matrix at $y=\Suc z$.

(ii) Rule \rulename{DE} opens
$\Q{y}\bigl(\sel^{+}(y,b,s)=\sel^{+}(y,c,t)\bigr)$; the opened parameter does
not occur. Its zero instance is $\Suc b=\Suc c$, which follows from $b=c$ by
congruence. Rule \rulename{QE} at the term $1$ gives
$\sel^{+}(1,b,s)=\sel^{+}(1,c,t)$, that is $s=t$.
\end{proof}

\begin{lemma}[Contextual Boolean cases]\label{seed:lem:succ-cases}
Let $b$ be Boolean in $\Gamma$ and let $E$ be a literal equation. If
$\Gamma,b=0\proves E$ and $\Gamma,b=1\proves E$, then $\Gamma\proves E$. The
context and $E$ may mention $b$.
\end{lemma}

\begin{proof}
Let $\Gamma\proves b=Z(r)$ and let $w$ be a variable not in $\Gamma$, $b$,
$r$, $E$. Apply \textsc{nc} to the native formula
$B(w)\defeq\mathcal I(b=Z(w);\,E)$. Its zero premise is
\cref{seed:lem:succ-imp}(i) applied to $\Gamma,b=Z(0)\proves E$, which is the
second given branch because $Z(0)=1$. Its successor premise is
\cref{seed:lem:succ-imp}(i) applied to $\Gamma,b=Z(\Suc z)\proves E$, which
is the first given branch because $Z(\Suc z)=0$. Rule \textsc{nc} gives
$\Gamma\proves B(r)$, and \cref{seed:lem:succ-imp}(ii) with
$\Gamma\proves b=Z(r)$ gives $E$. The context does not change during the
argument; the variable that varies is the fresh $w$.
\end{proof}

A branch whose hypotheses yield $0=\Suc 0$ is closed by the last statement of
\cref{seed:lem:succ-basic}.

\begin{remark}\label{seed:rem:succ-cases-variants}
\Cref{seed:lem:succ-decoder,seed:lem:succ-imp,seed:lem:succ-cases} are the
native bootstrap of \pkg{transient\_seed}.
\pkg{successor\_seed\_completion} proves the same case rule for terms with
values among $0,1,2$: a term $b$ has the \emph{functional range rule} for a
finite list $V$ of numerals if $\Gamma\proves K(r)$ for all $r\in V$ implies
$\Gamma\proves K(b)$; successor is cancelled on such terms by applying the
three-cycle $M$ of \cref{seed:lem:succ-cycle} and then $j$ twice, using
$M(w)=w$ and $j(j(j(w)))=w$ on $0,1,2$; and the guard formula is
$\D{a}{y}\bigl(\sel^{+}(y,x,u)=\sel^{+}(y,b,v)\bigr)$ evaluated at $x=r$ for
each $r\in V$. The predecessor proof below needs the two-valued case only.
\pkg{scalar\_successor} states its finite case analyses as abbreviations for
repeated \textsc{nc} and substitution; inside an induction step they are
instances of \cref{seed:lem:succ-cases}.
\end{remark}
% ed.: in scalar_successor the step of the ramp invariant M(F(n)) = N(n) and
% of the scaling identity split on N(n) with the induction hypothesis in the
% context; rule NC alone does not justify this, as transient_seed and
% successor_seed_completion both point out. The contextual case lemma closes
% the gap.

\subsection{The predecessor}

All statements of this subsection are derivable in $\mathsf{N}_R$, by
computation, the equality laws, \textsc{nc}, \textsc{ei} for the displayed
literal equation with its other parameters fixed, and
\cref{seed:lem:succ-cases}. They are therefore derivable in $\mathsf{T}_R$.
No step cancels a successor or uses a selector with unshifted branches.
The plan is that of \cref{base:thm:pred}: parity, a three-cycle, a ramp, a
factor map, a comparator and a saturating counter. The steps that
\cref{base:thm:pred} takes by equational interpolation and the two-point rule
(\cref{base:lem:interp}(iv), (v)) need the selector $\sel$ with unshifted
branches; here they are replaced by the contextual cases of
\cref{seed:lem:succ-cases}, which split a Boolean value under the induction
hypothesis, and every term is adjusted to the seed $\Suc a$. The conjugacy
and non-saturation inductions of \cref{seed:lem:succ-sat} and the last
comparator equation have the shape of the corresponding steps of
\cref{base:thm:pred,base:lem:cmp}, with the shifted values.

\begin{lemma}[Shifted sum]\label{seed:lem:succ-oplus}
Let $a\oplus b\defeq R(b,a;u.\Suc u)$. Then $a\oplus0=\Suc a$,
$a\oplus\Suc b=\Suc(a\oplus b)$, and
\[
 0\oplus a=\Suc a,\qquad \Suc a\oplus b=\Suc(a\oplus b),\qquad
 a\oplus b=b\oplus a,\qquad (a\oplus b)\oplus c=a\oplus(b\oplus c).
\]
\end{lemma}

\begin{proof}
The first two equations are \eqref{seed:eq:R0} and \eqref{seed:eq:RS}. For
$0\oplus a=\Suc a$ use \textsc{ei} on $a$: $0\oplus0=1$ and
$0\oplus\Suc a=\Suc(0\oplus a)=\Suc\Suc a$. For
$\Suc a\oplus b=\Suc(a\oplus b)$ use \textsc{ei} on $b$:
$\Suc a\oplus0=\Suc\Suc a=\Suc(a\oplus0)$ and
$\Suc a\oplus\Suc b=\Suc(\Suc a\oplus b)=\Suc\Suc(a\oplus b)
=\Suc(a\oplus\Suc b)$. For commutativity use \textsc{ei} on $b$:
$a\oplus0=\Suc a=0\oplus a$ and
$a\oplus\Suc b=\Suc(a\oplus b)=\Suc(b\oplus a)=\Suc b\oplus a$. For
associativity use \textsc{ei} on $c$: both sides are $\Suc(a\oplus b)$ at
$0$, and
$(a\oplus b)\oplus\Suc c=\Suc((a\oplus b)\oplus c)=\Suc(a\oplus(b\oplus c))
=a\oplus(b\oplus\Suc c)$.
\end{proof}

In the standard model $a\oplus b=a+b+1$. Define the affine terms
\[
 A_2(a)\defeq\Suc^2(a\oplus a),\qquad
 A_3(a)\defeq\Suc^4((a\oplus a)\oplus a),\qquad
 A_2^{\circ}(a)\defeq\Suc^3(a\oplus a),
\]
with standard values $2a+3$, $3a+6$ and $2a+4$. \Cref{seed:lem:succ-oplus}
gives
\begin{equation}\label{seed:eq:affine}
 A_2(0)=3,\qquad A_2(\Suc a)=\Suc^2A_2(a),\qquad
 A_3(0)=6,\qquad A_3(\Suc a)=\Suc^3A_3(a).
\end{equation}

\begin{lemma}[A three-cycle]\label{seed:lem:succ-cycle}
Let
\[
 j(u)\defeq\sel^{+}\bigl(u,0,\sel^{+}(Z(N(u)),1,0)\bigr),\qquad
 X(n)\defeq R(n,0;u.j(u)),\qquad M(n)\defeq j(j(X(n))).
\]
Then $j(0)=1$, $j(1)=2$, $j(2)=0$, $M(0)=0$ and $M(\Suc n)=j(M(n))$.
\end{lemma}

\begin{proof}
By \eqref{seed:eq:selplus}, $j(0)=\Suc 0$ and
$j(\Suc z)=\sel^{+}(Z(N(\Suc z)),1,0)$. Since $N(1)=1$ and $N(2)=N(0)=0$,
this gives $j(1)=\sel^{+}(0,1,0)=2$ and $j(2)=\sel^{+}(1,1,0)=0$. Further
$X(0)=1$ and $X(\Suc n)=j(X(n))$, so $M(0)=j(j(1))=0$ and
$M(\Suc n)=j(j(j(X(n))))=j(M(n))$.
\end{proof}

\begin{lemma}[The ramp]\label{seed:lem:succ-ramp}
Let $G(n)\defeq R(n,0;u.\,u\oplus M(\Suc^2u))$. Then
\begin{gather}
 M(G(n))=\Suc(N(n)), \label{seed:eq:ramp-inv}\\
 N(n)=0\proves G(\Suc n)=\Suc G(n),\qquad
 N(n)=1\proves G(\Suc n)=\Suc^2G(n), \label{seed:eq:ramp-cases}\\
 G(\Suc^2n)=\Suc^3G(n), \label{seed:eq:ramp-two}\\
 N(n)=0\proves G(\Suc^3n)=\Suc^4G(n),\qquad
 N(n)=1\proves G(\Suc^3n)=\Suc^5G(n). \label{seed:eq:ramp-three}
\end{gather}
\end{lemma}

\begin{proof}
By the computation equations, $G(0)=1$ and
$G(\Suc n)=G(n)\oplus M(\Suc^2G(n))=G(n)\oplus j(j(M(G(n))))$.

Equation \eqref{seed:eq:ramp-inv} is proved by \textsc{ei} on $n$. The base
is $M(1)=j(M(0))=j(0)=1=\Suc(N(0))$. For the step keep the induction
hypothesis $M(G(n))=\Suc(N(n))$ in the context and apply
\cref{seed:lem:succ-cases} to the Boolean term $N(n)$. If $N(n)=0$, then
$M(G(n))=1$, $j(j(1))=0$, so $G(\Suc n)=G(n)\oplus0=\Suc G(n)$ and
$M(G(\Suc n))=j(M(G(n)))=j(1)=2=\Suc(Z(0))=\Suc(N(\Suc n))$. If $N(n)=1$,
then $M(G(n))=2$, $j(j(2))=1$, so
$G(\Suc n)=G(n)\oplus1=\Suc^2G(n)$ and
$M(G(\Suc n))=j(j(M(G(n))))=j(j(2))=1=\Suc(Z(1))=\Suc(N(\Suc n))$.

The two computations of $G(\Suc n)$ just made, now with
\eqref{seed:eq:ramp-inv} as a theorem, are \eqref{seed:eq:ramp-cases}. For
\eqref{seed:eq:ramp-two} apply \cref{seed:lem:succ-cases} to $N(n)$: if
$N(n)=0$ then $N(\Suc n)=1$ and
$G(\Suc^2n)=\Suc^2G(\Suc n)=\Suc^3G(n)$; if $N(n)=1$ then $N(\Suc n)=0$ and
$G(\Suc^2n)=\Suc G(\Suc n)=\Suc^3G(n)$. Finally
$G(\Suc^3n)=\Suc^3G(\Suc n)$ by \eqref{seed:eq:ramp-two}, and
\eqref{seed:eq:ramp-cases} gives \eqref{seed:eq:ramp-three}.
\end{proof}

From $G(0)=1$, $N(0)=0$ and \eqref{seed:eq:ramp-cases},
\eqref{seed:eq:ramp-two} one computes
$G(1)=2$, $G(2)=4$, $G(3)=5$, $G(4)=7$, $G(6)=10$. In the standard model
$G(n)=\lfloor 3n/2\rfloor+1$.

\begin{lemma}[Affine identities]\label{seed:lem:succ-affine}
\begin{gather*}
 N(A_2(a))=1,\qquad N(A_3(a))=N(a),\qquad \Suc G(A_2(a))=A_3(a),\\
 G(A_3(a))=A_3(G(a))\oplus N(a).
\end{gather*}
\end{lemma}

\begin{proof}
Each equation is proved by \textsc{ei} on $a$ with \eqref{seed:eq:affine}.

$N(A_2(0))=N(3)=N(1)=1$, and
$N(A_2(\Suc a))=N(\Suc^2A_2(a))=N(A_2(a))$.

$N(A_3(0))=N(6)=N(0)$, and
$N(A_3(\Suc a))=N(\Suc^3A_3(a))=Z(N(A_3(a)))=Z(N(a))=N(\Suc a)$.

$\Suc G(A_2(0))=\Suc G(3)=6=A_3(0)$, and
$\Suc G(A_2(\Suc a))=\Suc G(\Suc^2A_2(a))=\Suc^3\Suc G(A_2(a))
=\Suc^3A_3(a)=A_3(\Suc a)$ by \eqref{seed:eq:ramp-two}.

For the last equation the base is $G(A_3(0))=G(6)=10$ and
$A_3(G(0))\oplus N(0)=A_3(1)\oplus0=\Suc 9=10$. For the step keep the
induction hypothesis $G(A_3(a))=A_3(G(a))\oplus N(a)$ and apply
\cref{seed:lem:succ-cases} to $N(a)$; the left side of the goal is
$G(\Suc^3A_3(a))$ and $N(A_3(a))=N(a)$.

If $N(a)=0$, then by \eqref{seed:eq:ramp-three} and the hypothesis
\[
 G(\Suc^3A_3(a))=\Suc^4G(A_3(a))=\Suc^4(A_3(G(a))\oplus0)=\Suc^5A_3(G(a)),
\]
while $G(\Suc a)=\Suc G(a)$, $N(\Suc a)=1$ and
\[
 A_3(G(\Suc a))\oplus N(\Suc a)=\Suc^3A_3(G(a))\oplus1=\Suc^5A_3(G(a)).
\]
If $N(a)=1$, then
\[
 G(\Suc^3A_3(a))=\Suc^5G(A_3(a))=\Suc^5(A_3(G(a))\oplus1)=\Suc^7A_3(G(a)),
\]
while $G(\Suc a)=\Suc^2G(a)$, $N(\Suc a)=0$ and
\[
 A_3(G(\Suc a))\oplus N(\Suc a)=\Suc^6A_3(G(a))\oplus0=\Suc^7A_3(G(a)).
 \qedhere
\]
\end{proof}

\begin{lemma}[The factor map]\label{seed:lem:succ-factor}
Let $H(u)\defeq\sel^{+}(Z(N(u)),G(u),u)$. Then
\begin{gather*}
 N(u)=0\proves H(u)=u,\qquad N(u)=1\proves H(u)=\Suc G(u),\\
 H(2)=2,\qquad H(A_2(a))=A_3(a),\qquad H(A_3(a))=A_3(H(a)).
\end{gather*}
\end{lemma}

\begin{proof}
The first line is \eqref{seed:eq:selplus} with $Z(0)=1$ and $Z(1)=0$. Since
$N(2)=0$, $H(2)=2$. Since $N(A_2(a))=1$,
$H(A_2(a))=\Suc G(A_2(a))=A_3(a)$. For the last equation apply
\cref{seed:lem:succ-cases} to $N(a)$, using $N(A_3(a))=N(a)$. If $N(a)=0$,
both sides are $A_3(a)$. If $N(a)=1$, then by \cref{seed:lem:succ-affine}
\[
 H(A_3(a))=\Suc G(A_3(a))=\Suc(A_3(G(a))\oplus1)=\Suc^3A_3(G(a))
 =A_3(\Suc G(a))=A_3(H(a)). \qedhere
\]
\end{proof}

On standard numbers $H$ fixes even inputs and sends an odd $u$ to
$3(u+1)/2$. In the coordinate $u+3$ the terms $A_2$ and $A_3$ are doubling
and tripling, and an active step of $H$ multiplies by $3/2$. This explains
the comparator below; the proof does not use it.

\begin{lemma}[A comparator]\label{seed:lem:succ-compare}
Let
\[
 \alpha(x)\defeq R(x,0;u.A_2^{\circ}(u)),\qquad e(x)\defeq\Suc\alpha(x),
 \qquad \Phi(x,y)\defeq R(y,\alpha(x);u.H(u)),\qquad
 L(x,y)\defeq Z(N(\Phi(x,y))).
\]
Then $L(x,y)$ is Boolean, and
\[
 L(0,y)=1,\qquad L(\Suc x,0)=0,\qquad L(\Suc x,\Suc y)=L(x,y),\qquad
 L(\Suc n,n)=0 .
\]
\end{lemma}

\begin{proof}
By computation and \cref{seed:lem:succ-oplus}, $e(0)=\Suc\alpha(0)=2$ and
\[
 e(\Suc x)=\Suc A_2^{\circ}(\alpha(x))=\Suc^4(\alpha(x)\oplus\alpha(x))
 =\Suc^2(\Suc\alpha(x)\oplus\Suc\alpha(x))=A_2(e(x)).
\]
Also $\Phi(x,0)=\Suc\alpha(x)=e(x)$ by \eqref{seed:eq:R0}, and
$\Phi(x,\Suc y)=H(\Phi(x,y))$. The initializer of $\Phi$ is chosen so that
the successor-seeded initial value is $e(x)$; no predecessor is needed to
start the orbit there.

By \textsc{ei} on $y$, $\Phi(0,y)=2$: the base is $e(0)=2$ and the step is
$H(2)=2$. Hence $L(0,y)=Z(N(2))=Z(0)=1$. Next
$\Phi(\Suc x,0)=A_2(e(x))$ has $N$-value $1$, so $L(\Suc x,0)=Z(1)=0$.

By \textsc{ei} on $y$ with $x$ fixed,
$\Phi(\Suc x,\Suc y)=A_3(\Phi(x,y))$. At $y=0$ the left side is
$H(e(\Suc x))=H(A_2(e(x)))=A_3(e(x))=A_3(\Phi(x,0))$. In the step both
orbits apply $H$, and $H(A_3(v))=A_3(H(v))$ turns
$H(A_3(\Phi(x,y)))$ into $A_3(\Phi(x,\Suc y))$. Applying $Z\circ N$ and
$N(A_3(v))=N(v)$ gives $L(\Suc x,\Suc y)=L(x,y)$.

Finally $L(\Suc n,n)=0$ by \textsc{ei} on $n$ from the second and third
equations, as in \cref{base:lem:cmp}.
\end{proof}

The standard value of $L(x,y)$ is $1$ if $x\leq y$ and $0$ otherwise.

\begin{lemma}[Saturation]\label{seed:lem:succ-sat}
Let
\[
 K(c,u)\defeq\sel^{+}(c,u,u),\qquad
 \beta_m(u)\defeq K(Z(L(\Suc u,m)),u),\qquad
 V(t,m)\defeq R(t,0;u.\beta_m(u)).
\]
Then
\[
 V(\Suc t,\Suc^2k)=\Suc V(t,\Suc k),\qquad V(n,\Suc n)=\Suc n,\qquad
 V(\Suc n,\Suc n)=V(\Suc^2n,\Suc n)=\Suc n .
\]
\end{lemma}

\begin{proof}
$K(0,u)=\Suc u$ and $K(\Suc v,u)=u$, so \textsc{nc} on $c$ gives
$K(c,\Suc u)=\Suc K(c,u)$. With the comparator equations,
\[
 \beta_{\Suc m}(\Suc u)=K(Z(L(\Suc u,m)),\Suc u)=\Suc\beta_m(u),\qquad
 \beta_{\Suc^2k}(1)=K(Z(L(0,k)),1)=K(0,1)=2 .
\]
The first equation of the lemma is \textsc{ei} on $t$ with $k$ fixed. At
$t=0$, $V(1,\Suc^2k)=\beta_{\Suc^2k}(V(0,\Suc^2k))=\beta_{\Suc^2k}(1)=2
=\Suc V(0,\Suc k)$. In the step,
$V(\Suc^2t,\Suc^2k)=\beta_{\Suc^2k}(\Suc V(t,\Suc k))
=\Suc\beta_{\Suc k}(V(t,\Suc k))=\Suc V(\Suc t,\Suc k)$.

The second equation is a separate \textsc{ei} on $n$: $V(0,1)=1$, and
$V(\Suc n,\Suc^2n)=\Suc V(n,\Suc n)$ by the first equation at $t=k=n$, an
instance of a completed lemma. For the third,
\[
 \beta_{\Suc n}(\Suc n)=K(Z(L(\Suc^2n,\Suc n)),\Suc n)
 =K(Z(L(\Suc n,n)),\Suc n)=K(1,\Suc n)=\Suc n,
\]
so one and then two further steps of the recursor from $V(n,\Suc n)=\Suc n$
leave the value unchanged.
\end{proof}

\begin{lemma}[A capped predecessor]\label{seed:lem:succ-capped}
Let $Q(m)\defeq R\bigl(m,0;u.K(Z(L(\Suc^2u,m)),u)\bigr)$. Then
$Q(0)=Q(1)=1$ and $Q(\Suc^2k)=\Suc k$.
\end{lemma}

\begin{proof}
$Q(0)=1$ by \eqref{seed:eq:R0}, and
$Q(1)=K(Z(L(3,1)),1)=K(Z(L(2,0)),1)=K(1,1)=1$. For $m=\Suc^2k$ the
comparator equations give
$L(\Suc^2u,\Suc^2k)=L(u,k)=L(\Suc u,\Suc k)$, so the body of $Q$ at this
parameter equals $\beta_{\Suc k}(u)$ for a generic $u$. By
\cref{seed:lem:succ-unique}, at an independent counter $t$ with $k$ fixed,
\[
 R\bigl(t,0;u.K(Z(L(\Suc^2u,\Suc^2k)),u)\bigr)=V(t,\Suc k).
\]
Substituting $t=\Suc^2k$ and \cref{seed:lem:succ-sat} give
$Q(\Suc^2k)=V(\Suc^2k,\Suc k)=\Suc k$.
\end{proof}

\begin{theorem}[Predecessor]\label{seed:thm:succ-pred}
Let
\[
 P(m)\defeq\sel^{+}\bigl(Z(L(3,m)),\;Q(Q(m)),\;L(2,m)\bigr).
\]
$\mathsf{N}_R$ and $\mathsf{T}_R$ derive
\[
 P(0)=0,\qquad P(\Suc n)=n .
\]
Consequently both derive $\Suc x=\Suc y\proves x=y$.
\end{theorem}

\begin{proof}
For $m\in\{0,1,2\}$ the comparator equations give $L(3,m)=0$, so
$Z(L(3,m))=1$ and $P(m)=L(2,m)$ by \eqref{seed:eq:selplus}; and
$L(2,0)=0$, $L(2,1)=L(1,0)=0$, $L(2,2)=L(0,0)=1$. Hence $P(0)=0$,
$P(1)=0$ and $P(2)=1$. For $m=\Suc^3k$, $L(3,m)=L(0,k)=1$, so the zero branch
of $\sel^{+}$ applies and, by \cref{seed:lem:succ-capped},
\[
 P(\Suc^3k)=\Suc Q(Q(\Suc^3k))=\Suc Q(\Suc^2k)=\Suc^2k .
\]
Rule \textsc{nc} on $n$, and in its successor case \textsc{nc} once more,
give $P(\Suc n)=n$ from $P(1)=0$, $P(2)=1$ and $P(\Suc^3k)=\Suc^2k$.
Injectivity follows by applying $P$ to $\Suc x=\Suc y$.
\end{proof}

\begin{remark}[The three constructions]\label{seed:rem:succ-three}
The construction above is that of \pkg{successor\_seed\_completion}: the
orbit $G$ is started at the seed $1$ that \eqref{seed:eq:R0} provides, the
affine terms absorb the shift, and the capped predecessor $Q$ is corrected at
three small arguments. The two other packages reach \cref{seed:thm:succ-pred}
by different terms.

\pkg{scalar\_successor} first builds a term $\tau$ with $\tau(0)=1$,
$\tau(1)=0$, $\tau(\Suc^2n)=\Suc^2n$ and $\tau(\tau(n))=n$, from the parity
term $U$, a Boolean conjunction $Z(\sel^{+}(a,0,Z(b)))$ and two finite orbits.
Conjugation gives an iterator that starts at zero,
\[
 K_0(n;u.h)\defeq\tau\bigl(R(n,0;v.\tau(h\subst{\tau(v)}{u}))\bigr),\qquad
 K_0(0;u.h)=0,\quad K_0(\Suc n;u.h)=h\subst{K_0(n;u.h)}{u},
\]
for every body $h$. With the shifted sum $a\oplus b$, the terms
$a\oplus a$ and $(a\oplus a)\oplus a$, a ramp $F(n)\defeq
K_0(n;u.\,u\oplus M'(u))$ over a three-state orbit $M'$ with the invariant
$M'(F(n))=N(n)$, and a factor map, it obtains a strict comparator $V$ with
$V(x,0)=0$, $V(0,\Suc y)=1$, $V(\Suc x,\Suc y)=V(x,y)$. The predecessor is
then the saturated zero-start orbit
$P(m)\defeq K_0\bigl(m;u.\sel^{+}(Z(V(\Suc u,m)),u,u)\bigr)$, with no
correction.

\pkg{transient\_seed} works in $\mathsf{T}_R$ and observes that
$f(R(n,a;u.f))$ starts at $f(\Suc a)$: a seed $\Suc a$ outside an intended
orbit is used once to enter it. Doubling and tripling orbits
$d_2(\Suc n)=\Suc^2d_2(n)$, $d_3(\Suc n)=\Suc^3d_3(n)$ that start at zero are
obtained in this way, then a ramp and a factor map, and a comparator $L$ with
the equations of \cref{seed:lem:succ-compare} and the persistence law
$L(x,y)=1\to L(x,\Suc y)=1$. The predecessor is a reset cycle: with
$F_m(u)\defeq\sel^{+}(Z(L(\Suc u,m)),u,0)$, which increments below $m$ and
resets to $0$ from $m$ on, and
\[
 J_m(t)\defeq F_m\bigl(R(t,m;u.F_m(u))\bigr),\qquad
 K_m(t)\defeq F_m\bigl(R(t,m;u.F_m(F_m(u)))\bigr),
\]
one has $J_m(0)=K_m(0)=0$ because $F_m(\Suc m)=0$. Induction on the literal
equation $\sel^{+}(L(t,m),0,J_m(t))=\sel^{+}(L(t,m),0,t)$ gives
$L(t,m)=1\to J_m(t)=t$; two further inductions give the period
$J_m(m\oplus t)=J_m(t)$ and $K_m(\Suc a)=J_m(a\oplus\Suc a)$; and
$P(n)\defeq K_n(n)$ satisfies $P(\Suc a)=J_{\Suc a}(a)=a$. The package then
transfers the proof to $\mathsf{N}_R$ by
\cref{seed:lem:succ-cases}.

The three proofs use the same resources: literal-equation induction with
fixed parameters, numerical cases, and Boolean cases in a context.
\end{remark}

\subsection{Recovery of the iterator}

\begin{theorem}[Initialized iteration]\label{seed:thm:succ-iter}
For every body $h$ and a fresh variable $v$ put
\[
 \widehat I(n,a;u.h)\defeq
 P\bigl(R(n,a;\,v.\,\Suc(h\subst{P(v)}{u}))\bigr).
\]
$\mathsf{N}_R$ and $\mathsf{T}_R$ derive
\[
 \widehat I(0,a;u.h)=a,\qquad
 \widehat I(\Suc n,a;u.h)=h\subst{\widehat I(n,a;u.h)}{u}.
\]
In particular $\widehat I(c,a;u.b)$ with $u\notin\FV(b)$ is a selector with
the unshifted values $a$ and $b$.
\end{theorem}

\begin{proof}
Write $r(n)\defeq R(n,a;v.\Suc(h\subst{P(v)}{u}))$. By \eqref{seed:eq:R0}
and \cref{seed:thm:succ-pred}, $\widehat I(0,a;u.h)=P(\Suc a)=a$. By
\eqref{seed:eq:RS},
$r(\Suc n)=\Suc(h\subst{P(r(n))}{u})$, so
$\widehat I(\Suc n,a;u.h)=P(\Suc(h\subst{P(r(n))}{u}))
=h\subst{P(r(n))}{u}=h\subst{\widehat I(n,a;u.h)}{u}$. The equation
$P(\Suc x)=x$ is used at arbitrary terms; no invariant about the state
$r(n)$ is needed.
\end{proof}

\begin{remark}
The recursor, the candidate equations and the compiler $\widehat I$ are
stated in an earlier draft of the series, which leaves the term $P$ open; the
three packages supply it. In the opposite direction the predecessor is the
only obstacle: a successor-seeded recursor whose update also receives the old
counter has an immediate predecessor (see \cref{seed:lem:tied-macros} for the
analogous computation).
\end{remark}

\subsection{Comparison of the two term languages}

\begin{definition}\label{seed:def:succ-trans}
The translation $t\mapsto t^{\iota}$ from $R$-terms to the terms of
\cref{base:sec:terms} and the translation $t\mapsto t^{\rho}$ in the opposite
direction are defined by
\begin{align*}
 x^{\iota}&\defeq x, & 0^{\iota}&\defeq0, &
 R(n,a;u.h)^{\iota}&\defeq\It{n^{\iota}}{\Suc(a^{\iota})}{u}{h^{\iota}},\\
 x^{\rho}&\defeq x, & 0^{\rho}&\defeq0, &
 (\Suc t)^{\rho}&\defeq\Suc(t^{\rho}),\qquad
 \It{n}{a}{u}{h}^{\rho}\defeq\widehat I(n^{\rho},a^{\rho};u.h^{\rho}),
\end{align*}
where $\Suc(t^{\rho})$ on the right is the abbreviation
$R(0,t^{\rho};u.0)$. Formulas, native or ordinary, are translated by
translating their terms. Both translations commute with substitution and
preserve free variables and the property of being a literal equation.
\end{definition}

Note that $(\Suc a)^{\iota}=\It{0}{\Suc(a^{\iota})}{u}{0}$, which is not
literally $\Suc(a^{\iota})$ but equals it by the zero computation equation of
the iterator.

\begin{theorem}[Equivalence of the term languages]\label{seed:thm:succ-terms}
\leavevmode
\begin{enumerate}[label=(\alph*),nosep]
\item If $\Gamma\proves A$ is derivable in $\mathsf{N}_R$, then
$\Gamma^{\iota}\proves A^{\iota}$ is derivable in $\Nat$. If
$\Gamma\proves A$ is derivable in $\mathsf{T}_R$, then
$\Gamma^{\iota}\proves A^{\iota}$ is derivable in $\Teq$.
\item If $\Gamma\proves A$ is derivable in $\Nat$, then
$\Gamma^{\rho}\proves A^{\rho}$ is derivable in $\mathsf{N}_R$. If
$\Gamma\proves A$ is derivable in $\Teq$, then
$\Gamma^{\rho}\proves A^{\rho}$ is derivable in $\mathsf{T}_R$.
\item $\mathsf{N}_R$ and $\mathsf{T}_R$ derive $t^{\iota\rho}=t$ for every
$R$-term $t$; $\Nat$ and $\Teq$ derive $t^{\rho\iota}=t$ for every iterator
term $t$. Hence $A^{\iota\rho}\approx A$ in $\mathsf{N}_R$ and
$A^{\rho\iota}\approx A$ in $\Nat$ for native formulas, and
$A^{\iota\rho}\leftrightarrow A$, $A^{\rho\iota}\leftrightarrow A$ are
provable in $\mathsf{T}_R$, $\Teq$ for ordinary formulas.
\item Consequently $\Gamma\proves A$ is derivable in $\mathsf{N}_R$ if and
only if $\Gamma^{\iota}\proves A^{\iota}$ is derivable in $\Nat$, and
likewise for $\mathsf{T}_R$ and $\Teq$, and for $(\cdot)^{\rho}$ in the other
direction. The contradictions correspond: $(0=\Suc 0)^{\iota}$ is
interderivable with $0=\Suc 0$, and $(0=\Suc 0)^{\rho}$ is $0=\Suc 0$.
\end{enumerate}
\end{theorem}

\begin{proof}
(a) By induction on the derivation. Each rule instance other than a
computation equation is translated into an instance of the same rule,
because the translation commutes with substitution, keeps literal equations
literal and keeps parameters fresh. One adjustment is needed where a rule
displays a successor: the premise of \rulename{DI} by cases translates to
$\Gamma^{\iota},B^{\iota}(t^{\iota},0)\proves
B^{\iota}(t^{\iota},\It{0}{\Suc z}{u}{0})$, and the premise of
\rulename{QI}, or of the induction rule of $\mathsf{T}_R$, to
$\Gamma^{\iota},E^{\iota}(z)\proves E^{\iota}(\It{0}{\Suc z}{u}{0})$.
Replacing $\It{0}{\Suc z}{u}{0}$ by $\Suc z$ along the zero computation
equation, by the Leibniz rule \cref{base:thm:leibniz}, gives the
premise of the same rule in $\Nat$ or $\Teq$. The translation of
\eqref{seed:eq:R0} is
$\It{0}{\Suc a^{\iota}}{u}{h^{\iota}}=\It{0}{\Suc a^{\iota}}{u}{0}$; both
sides equal $\Suc a^{\iota}$ by the zero computation equation. The
translation of \eqref{seed:eq:RS} is
\[
 \It{\It{0}{\Suc n^{\iota}}{u}{0}}{\Suc a^{\iota}}{u}{h^{\iota}}
 =h^{\iota}\subst{\It{n^{\iota}}{\Suc a^{\iota}}{u}{h^{\iota}}}{u},
\]
which follows from $\It{0}{\Suc n^{\iota}}{u}{0}=\Suc n^{\iota}$, congruence
and the successor computation equation. For $\mathsf{T}_R$ the logical rules
are translated identically; coverage translates to
$\forall x\,(x=0\lor\exists y\,x=\It{0}{\Suc y}{u}{0})$ and separation to
$0=\It{0}{\Suc 0}{u}{0}\to\bot$, both equivalent in $\Teq$ to the
corresponding axioms.

(b) By induction on the derivation. Here $(\Suc t)^{\rho}$ is literally
$\Suc(t^{\rho})$, so every rule instance is translated into an instance of
the same rule, and coverage, separation and induction are translated into
themselves. The two computation equations of the iterator are translated
into
$\widehat I(0,a^{\rho};u.h^{\rho})=a^{\rho}$ and
$\widehat I(\Suc n^{\rho},a^{\rho};u.h^{\rho})
=h^{\rho}\subst{\widehat I(n^{\rho},a^{\rho};u.h^{\rho})}{u}$, which are
instances of \cref{seed:thm:succ-iter}.

(c) First, $t^{\iota\rho}=t$ by induction on the $R$-term $t$. The cases of
variables and $0$ are trivial. For $t=R(n,a;u.h)$,
\[
 t^{\iota\rho}
 =P\Bigl(R\bigl(n^{\iota\rho},\Suc(a^{\iota\rho});\,
   v.\,\Suc(h^{\iota\rho}\subst{P(v)}{u})\bigr)\Bigr).
\]
By the induction hypothesis, congruence and
\cref{seed:lem:succ-unique} (for the body, whose induction hypothesis
$h^{\iota\rho}=h$ is derivable in the empty context with $u$ free), this
equals $P\bigl(R(n,\Suc a;v.\Suc(h\subst{P(v)}{u}))\bigr)$. The equation
\[
 P\bigl(R(z,\Suc a;\,v.\,\Suc(h\subst{P(v)}{u}))\bigr)=R(z,a;u.h)
\]
is proved by \textsc{ei} on $z$: at $0$ the left side is
$P(\Suc\Suc a)=\Suc a=R(0,a;u.h)$; at $\Suc z$ it is
$P(\Suc(h\subst{P(R(z,\Suc a;\ldots))}{u}))=h\subst{P(R(z,\Suc a;\ldots))}{u}$,
which by the induction hypothesis is
$h\subst{R(z,a;u.h)}{u}=R(\Suc z,a;u.h)$.

Second, $t^{\rho\iota}=t$ by induction on the iterator term $t$. By (a)
applied to \cref{seed:thm:succ-pred}, the term template
$P^{\iota}(m)\defeq(P(m))^{\iota}$ satisfies $P^{\iota}(0)=0$ and
$P^{\iota}(\It{0}{\Suc x}{u}{0})=x$, hence $P^{\iota}(\Suc x)=x$, in $\Nat$
and $\Teq$. For $t=\Suc s$,
$t^{\rho\iota}=\It{0}{\Suc(s^{\rho\iota})}{u}{0}=\Suc(s^{\rho\iota})=\Suc s$.
For $t=\It{n}{a}{u}{h}$,
\[
 t^{\rho\iota}=P^{\iota}\Bigl(\It{n^{\rho\iota}}{\Suc(a^{\rho\iota})}{v}
 {\It{0}{\Suc(h^{\rho\iota}\subst{P^{\iota}(v)}{u})}{w}{0}}\Bigr),
\]
which by the induction hypothesis, the zero computation equation and
uniqueness of iteration for the iterator (\cref{base:lem:uniq}, applied to
the iteration with the first body as the template $f$) equals
$P^{\iota}\bigl(\It{n}{\Suc a}{v}{\Suc(h\subst{P^{\iota}(v)}{u})}\bigr)$.
Induction for the literal equation
$P^{\iota}\bigl(\It{z}{\Suc a}{v}{\Suc(h\subst{P^{\iota}(v)}{u})}\bigr)
=\It{z}{a}{u}{h}$ proves the rest: at $0$ both sides are $a$, and at
$\Suc z$ the left side is
$h\subst{P^{\iota}(\It{z}{\Suc a}{v}{\ldots})}{u}$, which by the induction
hypothesis is $h\subst{\It{z}{a}{u}{h}}{u}=\It{\Suc z}{a}{u}{h}$.

The statements on formulas follow from
\cref{seed:lem:succ-leibniz}(i), in $\mathsf{N}_R$ and in $\Nat$, and from
replacement of provably equal terms under quantifiers in ordinary logic.

(d) If $\Gamma^{\iota}\proves A^{\iota}$ is derivable in $\Nat$, then
$\Gamma^{\iota\rho}\proves A^{\iota\rho}$ is derivable in $\mathsf{N}_R$ by
(b), and $\Gamma\proves A$ follows by (c) and cut. The other cases are the
same. For the contradiction, $(0=\Suc 0)^{\iota}$ is
$0=\It{0}{\Suc 0}{u}{0}$.
\end{proof}

\begin{corollary}\label{seed:thm:succ-bridge}
$\mathsf{N}_R$ and $\mathsf{T}_R$ are bi-interpretable with $\Nat$, $\Teq$
and $\iPRA$: the interpretations are the translations of
\cref{seed:def:succ-trans} composed with those of \cref{base:thm:bridge};
the domains are the whole domain, equality is actual equality, each term is
provably equal to its double translation, so that the graph $x=y$ is the
isomorphism of both numerical composites, and every formula is provably
equivalent to its double translation. The theories are equiconsistent.
In $\iPRA$ the recursor with body $h$ is the function defined by primitive
recursion with initial value $\Suc a$ and a step that ignores the counter.
\end{corollary}

\begin{proof}
\Cref{seed:thm:succ-terms} gives the statement for $\Nat$ and $\Teq$ with
the identity as numerical map. Composition with \cref{base:thm:bridge}
gives the rest; the description in $\iPRA$ is the composite of
$(\cdot)^{\iota}$ with the interpretation of the iterator.
\end{proof}

Consequently every result that \cref{base:thm:pred,base:thm:reflect,%
base:thm:primrec} state for $\Teq$ holds in $\mathsf{T}_R$ for the
translated terms: reflection of equality into a term, induction for
quantifier-free formulas, pairing, and a term with the defining equations
of each primitive recursive description. The predecessor of
\cref{base:thm:pred} is not needed: \cref{seed:thm:succ-pred} provides one
directly.

\begin{remark}[Rule forms in the packages]\label{seed:rem:succ-rules}
The three packages state the native calculus with two rule forms that differ
from \cref{base:sec:kernel}: transport is formula transport, from $s=t$ and
$F(s)$ infer $F(t)$ for every native formula template $F$; and progress
introduces the paired binder, from $\Gamma,E(t,z)\proves E(t,\Suc z)$ infer
$\Gamma\proves\D{x}{y}E(x,y)$ for a literal equation $E$. Both calculi derive
the same sequents. Formula transport is
\cref{seed:lem:succ-leibniz}(ii). The progress rule on $\D{x}{y}$ is the
witnessed introduction \cref{base:cor:combined}(i), derived from
\rulename{QI} and \cref{base:lem:dI}. Conversely \rulename{QI} is
derived from the progress rule on $\D{v}{y}E(y)$ with $v$ unused, followed by
\rulename{DE} and \rulename{Hyp}; and the transport rule of
\cref{base:sec:kernel} is a case of formula transport combined with symmetry.
\pkg{successor\_seed\_completion} and \pkg{scalar\_successor} also describe a
\emph{computation-supply} presentation used by their checkers: instead of a
computation leaf, a rule passes from $\Gamma,H\proves A$ to
$\Gamma\proves A$ for a computation equation $H$. The two presentations
derive the same sequents by cut in one direction and \rulename{Hyp} in the
other.
\end{remark}

%======================================================================
\section{Tied-seed recursors}
\label{seed:sec:tied}

This section draws on \pkg{bounded\_initializer},
\pkg{monus\_and\_neutral\_major} and the obstruction section of
\pkg{transient\_seed}.

\subsection{The recursor and its elementary terms}

\begin{definition}\label{seed:def:tied}
The tied-seed recursor $B(n;u,v.h)$ binds the distinct variables $u$ (old
counter) and $v$ (old state) in $h$. Its computation equations are
\begin{align}
 B(0;u,v.h)&=\Suc\bigl(h\subst{0}{u}\subst{0}{v}\bigr), \label{seed:eq:B0}\\
 B(\Suc n;u,v.h)&=h\subst{n}{u}\subst{B(n;u,v.h)}{v}. \label{seed:eq:BS}
\end{align}
The same body determines the value at $0$ and every later transition. Put
\[
 C(n,a)\defeq B(n;u,v.a),\quad W(n)\defeq B(n;u,v.u),\quad
 P(n)\defeq W(C(n,n)),\quad Z(n)\defeq C(n,0),
\]
with $u,v\notin\FV(a)$.
\end{definition}

\begin{lemma}\label{seed:lem:tied-macros}
From \eqref{seed:eq:B0}, \eqref{seed:eq:BS} and the equality laws alone,
\[
 C(0,a)=\Suc a,\quad C(\Suc n,a)=a,\quad W(0)=1,\quad W(\Suc n)=n,\quad
 P(0)=0,\quad P(\Suc n)=n,\quad Z(0)=1,\quad Z(\Suc n)=0 .
\]
Hence $\Suc x=\Suc y\proves x=y$, and $\Suc x=0\proves 0=\Suc 0$.
\end{lemma}

\begin{proof}
The first four equations are instances of the computation equations.
$P(0)=W(C(0,0))=W(1)=0$ and $P(\Suc n)=W(C(\Suc n,\Suc n))=W(\Suc n)=n$.
Applying $P$ to $\Suc x=\Suc y$ gives $x=y$. From $\Suc x=0$, applying $P$
gives $x=0$, and substitution gives $\Suc 0=0$.
\end{proof}

The successor can be removed from the primitive signature: $\sigma(a)\defeq
C(0,a)$ satisfies $\sigma(a)=\Suc a$, and writing $\sigma$ for $\Suc$ in
\eqref{seed:eq:B0}, \eqref{seed:eq:BS} and in every axiom gives an equivalent
presentation on the terms $x\mid0\mid B(n;u,v.h)$ and the other primitive
operations. The constant-body instance of \eqref{seed:eq:B0} becomes
tautological; \eqref{seed:eq:BS} for it and all other instances keep their
content. The section keeps $\Suc$ for readability.

\begin{lemma}[Gate and zero-start recursion]\label{seed:lem:tied-gate}
Assume induction for literal equations (rule \textsc{ei} of
\cref{seed:lem:succ-basic}). Let
\[
 G(n,a)\defeq P\bigl(P(B(a;u,v.C(Z(n),v)))\bigr),\qquad
 B_0(n;u,v.h)\defeq
 B\bigl(\Suc n;\,r,s.\,G(r,h\subst{P(r)}{u}\subst{s}{v})\bigr).
\]
Then
\[
 G(0,a)=0,\qquad G(\Suc n,a)=a,\qquad
 B_0(0;u,v.h)=0,\qquad
 B_0(\Suc n;u,v.h)=h\subst{n}{u}\subst{B_0(n;u,v.h)}{v}.
\]
\end{lemma}

\begin{proof}
For the control $0$ the body of the inner recursor is $C(1,v)=v$, its value
at $0$ is $\Suc(C(1,0))=1$, and \textsc{ei} on $a$ shows that the value at
$a$ is $1$; so $G(0,a)=P(P(1))=0$. For the control $\Suc n$ the body is
$C(0,v)=\Suc v$, the value at $0$ is $\Suc(\Suc 0)=2$, and \textsc{ei} on
$a$ shows that the value at $a$ is $\Suc\Suc a$; so
$G(\Suc n,a)=P(P(\Suc\Suc a))=a$. In both inductions $n$ is a fixed
parameter.

Let $k(r,s)\defeq G(r,h\subst{P(r)}{u}\subst{s}{v})$, so that $k(0,s)=0$
and $k(\Suc r,s)=h\subst{r}{u}\subst{s}{v}$. Then
$B(0;r,s.k)=\Suc(k(0,0))=1$, $B(1;r,s.k)=k(0,1)=0$, and
$B(\Suc\Suc n;r,s.k)=k(\Suc n,B(\Suc n;r,s.k))
=h\subst{n}{u}\subst{B(\Suc n;r,s.k)}{v}$. These are the equations for
$B_0$.
\end{proof}

The gate $G$ is a one-sided selector: it returns $0$ or its second argument.
The term $P(P(B(a;u,v.C(n,v))))$ is the complementary mask, with values $a$
at $n=0$ and $0$ at $n=\Suc m$.

\subsection{The one-coordinate bound}

The independent initializer cannot be recovered by a term. A function has a
\emph{bounding coordinate} in the sense of \cref{obs:def:bc}: some argument
$x_i$ and a nondecreasing unary primitive recursive $F$ with
$f(\vec x)\leq F(x_i)$; the general obstruction for tied seeds is
\cref{obs:cor:tied-seed}, and its two generalizations are
\cref{obs:thm:single-body,obs:thm:circuit}, both instances of the barrier for
adaptive query programs, \cref{obs:thm:adaptive}. The case of this chapter is the
following instance.

\begin{corollary}[One-coordinate bound]\label{seed:thm:tied-bound}
Consider the terms built from variables, $0$, $\Suc$, $B$, and optionally a
binary $\varepsilon$ with values in $\{0,1\}$ and a binary $d$ with
$d(x,y)\leq x$, in the standard model with \eqref{seed:eq:B0},
\eqref{seed:eq:BS}. The function denoted by a term is bounded by a constant
or has a bounding coordinate among its free variables, with a nondecreasing
primitive recursive bound that is obtained by recursion on the term.
\end{corollary}
% ed.: the packages prove the bound for the terms x | 0 | B (transient_seed)
% and with S and epsilon (bounded_initializer); monus_and_neutral_major
% asserts it for its terms with d. The case of d is added here.

\begin{proof}
The recursor $B$ is the constructor $\mathsf B$ of the tied-seed language of
\cref{obs:def:languages}(b), with the same computation equations
\eqref{obs:eq:B0}, \eqref{obs:eq:BS}. The successor is a unary primitive
recursive operation; $\varepsilon$ is globally bounded by $1$; and $d$ has
the bounding coordinate $x$ with the bound $F(z)=z$. By
\cref{obs:cor:tied-seed} every term denotes a function that is globally
bounded or has a bounding coordinate. The proof of that corollary, a
recursion on the term with the four cases of a constant body and of a
passive, counter or state coordinate of the body, composed with the bound of
the counter term by \cref{obs:lem:envelope}(iii), constructs the
nondecreasing primitive recursive bound.
\end{proof}

\begin{corollary}\label{seed:cor:tied-noterms}
No term of these languages denotes, in the standard model, any of the
following functions: the selector $\sel(n,a,b)$ with $\sel(0,a,b)=a$ and
$\sel(n+1,a,b)=b$; addition; an injection $\mathbb N^2\to\mathbb N$.
\end{corollary}

\begin{proof}
By \cref{seed:thm:tied-bound} every such function would be globally bounded
or have a bounding coordinate, and the argument of \cref{obs:cor:no-addition}
shows that none of the three has either property.
\end{proof}

\begin{remark}\label{seed:rem:tied-bound-scope}
The bound is in \pkg{transient\_seed} for the terms $x\mid0\mid B$, and in
\pkg{bounded\_initializer} with $\Suc$ and $\varepsilon$; the case of $d$ is
immediate from $d(x,y)\leq x$. The recursor $J(0;h)=1$,
$J(\Suc n;h)=h(n,J(n;h))$, which \pkg{scalar\_successor} mentions as the
reason for keeping the initializer in $R$, is a single-body constructor
scheme of \cref{obs:def:single-body} (with the constant seed combiner $1$),
so \cref{obs:thm:single-body} gives the same bound for it. The bound concerns the
functions denoted by terms. Sound logical rules cannot change it, so there is
no compiler from the iterator language to these terms. It does not show that
a theory on these terms fails to be bi-interpretable with $\iPRA$:
\cref{seed:thm:tied-biint} interprets each primitive recursive function by a
graph with finitely many value terms. The successor-seeded recursor
$R(n,a;u.h)$ escapes the bound because $a\oplus b$ is unbounded in each of
its two arguments. Related obstructions for tied seeds are in the package
\pkg{local\_scopes} and in \cref{obs:thm:single-body,obs:thm:circuit}.
\end{remark}

\subsection{The theory with a Boolean equality test}

\begin{definition}\label{seed:def:tbe}
$\mathsf{T}_{B\varepsilon}$ is ordinary intuitionistic first-order logic
with equality over the terms
$t\mathrel{::=}x\mid0\mid\Suc t\mid\varepsilon(t,t)\mid B(n;u,v.h)$, with
the axioms \eqref{seed:eq:B0}, \eqref{seed:eq:BS},
\begin{gather}
 \varepsilon(x,y)=0\lor\varepsilon(x,y)=1, \label{seed:eq:erange}\\
 \varepsilon(x,y)=1\leftrightarrow x=y, \label{seed:eq:ereflect}\\
 0=1\to\bot, \label{seed:eq:sep}
\end{gather}
universally closed, and induction for literal equations with side
assumptions. Coverage, injectivity of the successor and induction for
compound formulas are not axioms.
\end{definition}

The equality on the right of \eqref{seed:eq:ereflect} is the equality of the
logic on the whole domain. The axioms make numerical equality decidable and
give $\varepsilon(s,t)=0\leftrightarrow s\neq t$; they do not give excluded
middle for quantified formulas.

\begin{lemma}[Quantifier-free induction and coverage]\label{seed:lem:tbe-qf}
In $\mathsf{T}_{B\varepsilon}$ every quantifier-free formula $A$ has a term
$\chi_A$ with $\chi_A=0\lor\chi_A=1$ and $A\leftrightarrow\chi_A=1$.
Induction with side assumptions holds for quantifier-free formulas, and
$x=0\lor x=\Suc(P(x))$ is provable.
\end{lemma}

\begin{proof}
By recursion on $A$: $\chi_{s=t}\defeq\varepsilon(s,t)$, $\chi_\bot\defeq0$,
and on Boolean values $Z$ is negation and $G$ is conjunction by
\cref{seed:lem:tied-macros,seed:lem:tied-gate}; disjunction is
$Z(G(Z(b),Z(c)))$ and implication is $Z(G(b,Z(c)))$. The truth tables are
checked on the two values. The clause for implication is valid
intuitionistically because $A$ is decidable by its characteristic term.
Given a base and a step for $A(n)$, pass to the literal equation
$\chi_{A(n)}=1$, apply the induction rule and pass back. Apply the result to
$x=0\lor x=\Suc(P(x))$: the base is immediate, and the step holds by
$P(\Suc x)=x$ without the induction hypothesis.
\end{proof}

For a Boolean $b$ let $\operatorname{inc}(a,b)\defeq C(Z(b),a)$, so that
$\operatorname{inc}(a,0)=a$ and $\operatorname{inc}(a,1)=\Suc a$. Reflection
and \cref{seed:lem:tied-macros} give
\[
 \varepsilon(x,x)=1,\qquad\varepsilon(x,y)=\varepsilon(y,x),\qquad
 \varepsilon(\Suc x,\Suc y)=\varepsilon(x,y),\qquad\varepsilon(0,\Suc x)=0 .
\]
Define the \emph{hit term} and the order relations
\[
 \operatorname{hit}(n,y)\defeq
 B_0\bigl(n;r,v.\operatorname{inc}(v,\varepsilon(r,y))\bigr),\qquad
 x<y\iffdef\operatorname{hit}(y,x)=1,\qquad
 x\leq y\iffdef\operatorname{hit}(\Suc y,x)=1 .
\]
Thus $\operatorname{hit}(0,y)=0$ and
$\operatorname{hit}(\Suc n,y)
=\operatorname{inc}(\operatorname{hit}(n,y),\varepsilon(n,y))$.

\begin{lemma}[Hit identities]\label{seed:lem:tbe-hit}
$\mathsf{T}_{B\varepsilon}$ proves
\begin{gather*}
 \operatorname{hit}(\Suc n,0)=1,\quad
 \operatorname{hit}(\Suc n,\Suc y)=\operatorname{hit}(n,y),\quad
 \operatorname{hit}(n,n)=0,\quad\operatorname{hit}(\Suc n,n)=1,\\
 \operatorname{hit}(n,y)=0\lor\operatorname{hit}(n,y)=1,\qquad
 \operatorname{hit}(\Suc y,x)=Z(\operatorname{hit}(x,y)),\\
 \operatorname{hit}(n,y)=1\to\operatorname{hit}(\Suc n,y)=1,\qquad
 \operatorname{hit}(\Suc n,y)=1\to\operatorname{hit}(n,y)=1\lor n=y .
\end{gather*}
\end{lemma}

\begin{proof}
$\operatorname{hit}(\Suc n,0)=1$ by induction on $n$: the first step tests
$\varepsilon(0,0)=1$, and every later test $\varepsilon(\Suc n,0)$ is $0$.
$\operatorname{hit}(\Suc n,\Suc y)=\operatorname{hit}(n,y)$ by induction on
$n$ with $y$ fixed: the base is
$\operatorname{inc}(0,\varepsilon(0,\Suc y))=0$, and the step uses
$\varepsilon(\Suc n,\Suc y)=\varepsilon(n,y)$. Then
$\operatorname{hit}(n,n)=0$ by induction, and
$\operatorname{hit}(\Suc n,n)=\operatorname{inc}(0,1)=1$.

For the range use quantifier-free induction on $n$: if
$\varepsilon(n,y)=0$ the state is unchanged; if $\varepsilon(n,y)=1$ then
$n=y$, the old state is $\operatorname{hit}(n,n)=0$ and the new state is
$1$. If $\operatorname{hit}(n,y)=1$, then $n\neq y$ by the diagonal value,
so the next test is $0$ and the hit persists. The same case split gives the
last implication.

The duality $\operatorname{hit}(\Suc y,x)=Z(\operatorname{hit}(x,y))$ is
proved by induction on $x$ with $y$ fixed. At $x=0$ both sides are $1$. In
the step the left side is
$\operatorname{hit}(\Suc y,\Suc x)=\operatorname{hit}(y,x)$ and the right
side is $Z(\operatorname{inc}(\operatorname{hit}(x,y),\varepsilon(x,y)))$.
If $\varepsilon(x,y)=0$, the right side is $Z(\operatorname{hit}(x,y))$, and
the induction hypothesis reads
$\operatorname{inc}(\operatorname{hit}(y,x),0)=\operatorname{hit}(y,x)
=Z(\operatorname{hit}(x,y))$. If $\varepsilon(x,y)=1$, then $x=y$, the left
side is $0$ and the right side is $Z(\operatorname{hit}(\Suc x,x))=Z(1)=0$.
\end{proof}

\begin{lemma}[Order]\label{seed:lem:tbe-order}
In $\mathsf{T}_{B\varepsilon}$ the relations $<$ and $\leq$ form a discrete
linear order with least element $0$:
$x\leq y\leftrightarrow(x=y\lor x<y)$,
$x<\Suc y\leftrightarrow x\leq y$, exactly one of $x<y$, $x=y$, $y<x$
holds, both relations are transitive, the successor preserves and reflects
them, and $x\leq\Suc y$ with $x\neq\Suc y$ implies $x\leq y$.
\end{lemma}

\begin{proof}
The diagonal identities give $x\leq x$, $x\not<x$ and $x<\Suc x$. The
recurrence for $\operatorname{hit}(\Suc y,x)$, split on $\varepsilon(y,x)$,
gives the first equivalence. If $x\neq y$, then
$\operatorname{hit}(\Suc y,x)=\operatorname{hit}(y,x)$, and by duality
$\operatorname{hit}(x,y)$ is its Boolean complement; this is trichotomy. For
transitivity of $<$ keep $\operatorname{hit}(y,x)=1$ and induct on $z$ in
the quantifier-free formula
$\operatorname{hit}(z,y)=1\to\operatorname{hit}(z,x)=1$. The base has a
false antecedent. At $\Suc z$ the last implication of
\cref{seed:lem:tbe-hit} gives $\operatorname{hit}(z,y)=1$, and then the
induction hypothesis and persistence apply, or $z=y$, and then the retained
$\operatorname{hit}(y,x)=1$ and persistence apply. Transitivity of $\leq$
and the mixed laws follow by splitting equality from strict comparison, and
antisymmetry from trichotomy. The identity
$\operatorname{hit}(\Suc n,\Suc y)=\operatorname{hit}(n,y)$ gives
preservation and reflection by the successor.
$x<\Suc y\leftrightarrow x\leq y$ holds by definition, and the last claim
follows from it and the first equivalence. Finally
$\operatorname{hit}(\Suc x,0)=1$ is $0\leq x$, and $x\leq0$ gives $x=0$ by
trichotomy.
\end{proof}

Define the difference, the minimum, doubling and bounded addition by
\begin{gather*}
 \widehat d(x,y)\defeq
 B_0\bigl(x;u,v.\operatorname{inc}(v,\operatorname{hit}(\Suc u,y))\bigr),
 \qquad\operatorname{mn}(x,y)\defeq\widehat d(x,\widehat d(x,y)),\qquad
 \operatorname{db}(n)\defeq B_0(n;u,v.\Suc\Suc v),\\
 K_b(a)\defeq\widehat d(\operatorname{db}(b),\widehat d(b,a)),\qquad
 \operatorname{ad}_b(a,c)\defeq\widehat d(K_b(a),\widehat d(b,c)).
\end{gather*}
Thus $\widehat d(0,y)=0$ and
$\widehat d(\Suc x,y)
=\operatorname{inc}(\widehat d(x,y),\operatorname{hit}(\Suc x,y))$.

\begin{lemma}[Difference]\label{seed:lem:tbe-diff}
$\mathsf{T}_{B\varepsilon}$ proves
\begin{gather*}
 \widehat d(x,0)=x,\qquad\widehat d(\Suc x,\Suc y)=\widehat d(x,y),\qquad
 \widehat d(x,x)=0,\qquad\widehat d(x,y)=0\leftrightarrow x\leq y,\\
 \widehat d(x,\Suc y)=P(\widehat d(x,y)),\qquad\widehat d(x,y)\leq x,\\
 y\leq x\to\widehat d(\Suc x,y)=\Suc(\widehat d(x,y)),\qquad
 0<y\leq x\to\widehat d(x,P(y))=\Suc(\widehat d(x,y)),
\end{gather*}
and that $\widehat d$ is monotone in its first argument and antitone in its
second.
\end{lemma}

\begin{proof}
$\widehat d(x,0)=x$ by induction, since $\operatorname{hit}(\Suc x,0)=1$.
For $\widehat d(\Suc x,\Suc y)=\widehat d(x,y)$ induct on $x$ with $y$
fixed: the base uses $\operatorname{hit}(1,\Suc y)=0$, and the step uses
$\operatorname{hit}(\Suc\Suc x,\Suc y)=\operatorname{hit}(\Suc x,y)$.

Quantifier-free induction on $x$ with $y$ fixed proves
$\widehat d(x,y)=0\leftrightarrow\operatorname{hit}(x,y)=0$. In the step, if
$\operatorname{hit}(\Suc x,y)=1$ the new difference is a successor; if
$\operatorname{hit}(\Suc x,y)=0$, then $\operatorname{hit}(x,y)=0$ by
persistence, the old difference is $0$ and stays $0$. By duality
$\operatorname{hit}(x,y)=0$ is $\operatorname{hit}(\Suc y,x)=1$, which gives
the characterization of $\leq$ and $\widehat d(x,x)=0$.

For $\widehat d(x,\Suc y)=P(\widehat d(x,y))$ induct on $x$ with $y$ fixed.
In the step the left side is $\widehat d(x,y)$ and the right side is
$P(\operatorname{inc}(\widehat d(x,y),\operatorname{hit}(\Suc x,y)))$. If
the Boolean is $1$ the predecessor cancels the successor; if it is $0$ the
characterization gives $\widehat d(\Suc x,y)=0$, hence
$\widehat d(x,y)=0$ and both sides are $0$.

Quantifier-free induction on $x$ proves $\widehat d(x,y)\leq x$: the new
state is the old one or its successor. For monotonicity in the first
argument induct on $v$ in
$u\leq v\to\widehat d(u,z)\leq\widehat d(v,z)$ with $u$, $z$ fixed: at
$\Suc v$ either $u=\Suc v$, or $u\leq v$ by discreteness and the new state
is at least the old one; the base uses $u\leq0\to u=0$. For antitonicity keep
$y\leq z$ and induct on $x$ in $\widehat d(x,z)\leq\widehat d(x,y)$: if
$\operatorname{hit}(\Suc x,z)=0$ the left state stays and the right state
does not decrease; if it is $1$ then $z\leq x$, so $y\leq x$ and both states
increase.

The first of the two final implications is the defining recurrence, since
$y\leq x$ means $\operatorname{hit}(\Suc x,y)=1$. For the second,
$y=\Suc(P(y))$ by coverage, $P(y)<x$, so $\widehat d(x,P(y))\neq0$, and
coverage at this difference with
$\widehat d(x,\Suc(P(y)))=P(\widehat d(x,P(y)))$ gives the claim.
\end{proof}

\begin{lemma}[Minimum, doubling, bounded addition]\label{seed:lem:tbe-add}
$\mathsf{T}_{B\varepsilon}$ proves:
\begin{enumerate}[label=(\alph*),nosep]
\item $x\leq y\to\operatorname{mn}(x,y)=x$ and
$y\leq x\to\operatorname{mn}(x,y)=y$;
\item $\operatorname{db}$ is monotone, $n\leq\operatorname{db}(n)$ and
$\widehat d(\operatorname{db}(n),n)=n$;
\item $a\leq b\to\operatorname{ad}_b(a,0)=a$;
\item $a\leq b\land\Suc c\leq b\to
\operatorname{ad}_b(a,\Suc c)=\Suc(\operatorname{ad}_b(a,c))$;
\item $c\leq b\to\operatorname{ad}_b(0,c)=c$;
\item $\operatorname{ad}_b(b,b)=\operatorname{db}(b)$ and
$\operatorname{ad}_b(a,c)\leq\operatorname{db}(b)$ for all $a$, $c$;
\item $\operatorname{ad}_b$ is monotone in its second argument;
\item if $a,c\leq b$ and $a,c\leq b'$ then
$\operatorname{ad}_b(a,c)=\operatorname{ad}_{b'}(a,c)$.
\end{enumerate}
\end{lemma}

\begin{proof}
(a) The first claim uses $\widehat d(x,y)=0$ and $\widehat d(x,0)=x$. For the
second induct on $y$ in $y\leq x\to\widehat d(x,\widehat d(x,y))=y$ with $x$
fixed. The base is $\widehat d(x,x)=0$. If $\Suc y\leq x$, then $y<x$ and
$0<\widehat d(x,y)\leq x$, so by \cref{seed:lem:tbe-diff}
$\widehat d(x,\widehat d(x,\Suc y))=\widehat d(x,P(\widehat d(x,y)))
=\Suc(\widehat d(x,\widehat d(x,y)))=\Suc y$.

(b) $\operatorname{db}(0)=0$ and
$\operatorname{db}(\Suc n)=\Suc\Suc\operatorname{db}(n)$ give
$n\leq\operatorname{db}(n)$ by quantifier-free induction; monotonicity is an
induction on the larger argument that splits equality at the successor from
the discrete predecessor case. The identity is an induction on a literal
equation:
$\widehat d(\operatorname{db}(\Suc n),\Suc n)
=\widehat d(\Suc\operatorname{db}(n),n)
=\Suc(\widehat d(\operatorname{db}(n),n))=\Suc n$, using
$n\leq\operatorname{db}(n)$.

(c)--(h) Since $\widehat d(b,a)\leq b$, antitonicity and (b) give
$b=\widehat d(\operatorname{db}(b),b)\leq K_b(a)\leq\operatorname{db}(b)$,
and $K_b(0)=b$. If $\Suc a\leq b$, put $q\defeq\widehat d(b,a)$; then
$0<q\leq b\leq\operatorname{db}(b)$ and $\widehat d(b,\Suc a)=P(q)$, so
$K_b(\Suc a)=\widehat d(\operatorname{db}(b),P(q))=\Suc(K_b(a))$.
Quantifier-free induction on $a$ proves (c) under $a\leq b$: the base is
$\widehat d(b,b)=0$, and the step increments $\widehat d(K_b(a),b)$ by the
last identity and $b\leq K_b(a)$. For (d) put $q\defeq\widehat d(b,c)$; then
$0<q\leq b\leq K_b(a)$ and
$\operatorname{ad}_b(a,\Suc c)=\widehat d(K_b(a),P(q))
=\Suc(\widehat d(K_b(a),q))$. Statement (e) is a quantifier-free induction
on $c\leq b$ from (c) and (d). For (f) substitute $a=c=b$ and use
$\widehat d(b,b)=0$ twice; the bound is
$\widehat d(K_b(a),\widehat d(b,c))\leq K_b(a)\leq\operatorname{db}(b)$. For
(g), $c\leq c'$ gives $\widehat d(b,c')\leq\widehat d(b,c)$, and antitonicity
of the outer difference gives the claim. For (h) keep $a\leq b,b'$ and
induct on $c$ in
$c\leq\operatorname{mn}(b,b')\to
\operatorname{ad}_b(a,c)=\operatorname{ad}_{b'}(a,c)$: the base gives $a$ on
both sides, and (d) applies to both sides whenever
$\Suc c\leq\operatorname{mn}(b,b')$.
\end{proof}

Every induction in
\cref{seed:lem:tbe-hit,seed:lem:tbe-order,seed:lem:tbe-diff,seed:lem:tbe-add}
is on a quantifier-free formula with the other variables fixed. There is no
total binary addition term; $\operatorname{ad}_b$ adds below the bound $b$.

\begin{lemma}[Clipped initialized recursion]\label{seed:lem:tbe-clip}
Let
\[
 \sel_b(c,a,e)\defeq\operatorname{ad}_b\bigl(G(Z(c),a),G(c,e)\bigr),\qquad
 \operatorname{Rec}_b(n,a;u,v.h)\defeq
 B_0\Bigl(\Suc n;\,r,s.\,
 \sel_b\bigl(r,\operatorname{mn}(a,b),
 \operatorname{mn}(h\subst{P(r)}{u}\subst{s}{v},b)\bigr)\Bigr).
\]
If $a,e\leq b$, then $\sel_b(0,a,e)=a$ and $\sel_b(\Suc c,a,e)=e$. For all
arguments,
\[
 \operatorname{Rec}_b(0,a;u,v.h)=\operatorname{mn}(a,b),\qquad
 \operatorname{Rec}_b(\Suc n,a;u,v.h)
 =\operatorname{mn}\bigl(h\subst{n}{u}
 \subst{\operatorname{Rec}_b(n,a;u,v.h)}{v},\,b\bigr),
\]
and $\operatorname{Rec}_b(n,a;u,v.h)\leq b$.
\end{lemma}

\begin{proof}
By \cref{seed:lem:tied-gate} one of the two summands of $\sel_b$ is $0$ and
the other is $a$ at $c=0$ and $e$ at a successor;
\cref{seed:lem:tbe-add}(c) and (e) give the values. The equations for
$\operatorname{Rec}_b$ follow from those of $B_0$ and of $\sel_b$, since both
arguments $\operatorname{mn}(\cdot,b)$ are at most $b$; the bound is a
quantifier-free induction.
\end{proof}

The cap $b$ is a parameter that stays fixed during the recursion. Arbitrary
initialization is thus available once a cap has been fixed, in agreement
with \cref{seed:thm:tied-bound}: the value is bounded in the coordinate $b$.

\subsection{A compiler with caps}

Fix the primitive recursive descriptions built from zero, successor,
projections, composition and primitive recursion. The following construction
is a recursion on a description outside the theory; it produces terms and
proofs of $\mathsf{T}_{B\varepsilon}$ and adds no evaluator for codes of
descriptions.

\begin{theorem}[Bounded realization]\label{seed:thm:tbe-realize}
For every $k$-ary primitive recursive description $f$ there are terms
$\tau_f(b,x_1,\ldots,x_k)$ and $M_f(b)$ such that
$\mathsf{T}_{B\varepsilon}$ proves:
\begin{enumerate}[label=(\roman*),nosep]
\item $M_f$ is nondecreasing and $b\leq M_f(b)$;
\item $\tau_f(b,\vec x)\leq M_f(b)$ for all $b$, $\vec x$;
\item if $x_i\leq b$ and $x_i\leq c$ for all $i$, then
$\tau_f(b,\vec x)=\tau_f(c,\vec x)$.
\end{enumerate}
\end{theorem}

\begin{proof}
By recursion on $f$.

\emph{Initial descriptions.} $\tau_0\defeq0$,
$\tau_{\pi_i}(b,\vec x)\defeq\operatorname{mn}(x_i,b)$,
$\tau_{\Suc}(b,x)\defeq\Suc\operatorname{mn}(x,b)$, with $M(b)\defeq\Suc b$.
The properties follow from \cref{seed:lem:tbe-add}(a).

\emph{Composition.} Let $f(\vec x)=h(g_1(\vec x),\ldots,g_m(\vec x))$. Put
$Q\defeq M_{g_1}\circ\cdots\circ M_{g_m}$ (the identity if $m=0$),
$M_f\defeq M_h\circ Q$ and
\[
 \tau_f(b,\vec x)\defeq
 \tau_h\bigl(Q(b),\tau_{g_1}(b,\vec x),\ldots,\tau_{g_m}(b,\vec x)\bigr).
\]
Since the $M_{g_i}$ are nondecreasing and $z\leq M_{g_i}(z)$,
$M_{g_i}(b)\leq Q(b)$. (ii) is the bound for $\tau_h$. For (iii) the inner
values agree at the caps $b$ and $c$ by the hypothesis for the $g_i$, and
they are bounded by $Q(b)$ and by $Q(c)$; (iii) for $h$ gives the claim. (i)
holds by composition.

\emph{Primitive recursion.} Let $f(0,\vec x)=g(\vec x)$ and
$f(\Suc n,\vec x)=h(n,f(n,\vec x),\vec x)$. Put
\[
 F\defeq M_h\circ M_g\circ\Suc,\qquad
 I_F(t)\defeq B_0(t;u,v.F(v)),\qquad
 C_b\defeq I_F(\Suc\operatorname{db}(b)),\qquad M_f(b)\defeq C_b .
\]
$F$ is nondecreasing with $\Suc z\leq F(z)$ and $M_h(z)\leq F(z)$. Then
$I_F(0)=0$, $I_F(\Suc t)=F(I_F(t))$, $t\leq I_F(t)$ by induction, and $I_F$
is nondecreasing by a quantifier-free induction that uses discreteness of the
order; so $b\mapsto C_b$ is nondecreasing and $b\leq C_b$. For $j\leq b$
define the stage bound
\[
 \kappa_b(j)\defeq I_F\bigl(\Suc\operatorname{ad}_b(b,j)\bigr).
\]
\Cref{seed:lem:tbe-add} gives
\begin{equation}\label{seed:eq:stage}
 b\leq\kappa_b(0),\quad M_g(b)\leq\kappa_b(0),\qquad
 \Suc j\leq b\to\kappa_b(\Suc j)=F(\kappa_b(j)),\qquad
 j\leq b\to b\leq\kappa_b(j)\leq C_b .
\end{equation}
Indeed $\kappa_b(0)=I_F(\Suc b)=F(I_F(b))\geq F(b)\geq M_g(b)$; the
recurrence is \cref{seed:lem:tbe-add}(d); the upper bound is
\cref{seed:lem:tbe-add}(f) and the lower bound is monotonicity.

Write $n_b\defeq\operatorname{mn}(n,b)$ and
$\vec x_b\defeq(\operatorname{mn}(x_i,b))_i$, and define
\[
 \tau_f(b,n,\vec x)\defeq
 \operatorname{Rec}_{C_b}\bigl(n_b,\tau_g(b,\vec x_b);\,
 u,v.\,\tau_h(C_b,u,v,\vec x_b)\bigr).
\]
(ii) is the bound of \cref{seed:lem:tbe-clip}.

\emph{No clipping on bounded inputs.} Fix $\vec x\leq b$ and let
$r_b(j)\defeq\operatorname{Rec}_{C_b}(j,\tau_g(b,\vec x);
u,v.\tau_h(C_b,u,v,\vec x))$. Quantifier-free induction on $j$ proves
\begin{equation}\label{seed:eq:growth}
 j\leq b\to r_b(j)\leq\kappa_b(j).
\end{equation}
At $0$, $\tau_g(b,\vec x)\leq M_g(b)\leq\kappa_b(0)\leq C_b$, so the seed is
not clipped. Let $\Suc j\leq b$. The arguments $j$, $r_b(j)$, $\vec x$ of the
step are bounded by $\kappa_b(j)$ and by $C_b$, so (iii) for $h$ gives
\[
 \tau_h(C_b,j,r_b(j),\vec x)=\tau_h(\kappa_b(j),j,r_b(j),\vec x)
 \leq M_h(\kappa_b(j))\leq F(\kappa_b(j))=\kappa_b(\Suc j)\leq C_b .
\]
The value is not clipped, which proves \eqref{seed:eq:growth} and
\begin{equation}\label{seed:eq:unclipped}
 \Suc j\leq b\to r_b(\Suc j)=\tau_h(C_b,j,r_b(j),\vec x).
\end{equation}

\emph{Coherence.} Fix $b$, $c$ and $n,\vec x$ bounded by both. The minimum
identities remove the clipping of the inputs. Quantifier-free induction on
$j$ proves $(j\leq b\land j\leq c)\to r_b(j)=r_c(j)$: the base is (iii) for
$g$ and the absence of seed clipping; if $\Suc j\leq b,c$, the two states at
$j$ are equal and bounded by $C_b$ and $C_c$, as are the other arguments, so
(iii) for $h$ equates the two step values and \eqref{seed:eq:unclipped}
removes their clipping. At $j=n$ this is (iii) for $f$. The parameters
$b,c,n,\vec x$ are fixed and both induction formulas are quantifier-free.

Every body inserted into $B_0$ or $\operatorname{Rec}_b$ is the already
compiled term $\tau_h$ of the original language. No recursion over a
function introduced later by a graph is used.
\end{proof}

For a nonempty tuple $\vec x=(x_1,\ldots,x_k)$ let
\[
 C_i(\vec x)\defeq\bigwedge_{j<i}x_j<x_i\land\bigwedge_{j>i}x_j\leq x_i
\]
say that $i$ is the first largest coordinate. By \cref{seed:lem:tbe-order}
these $k$ quantifier-free cases are exhaustive and pairwise disjoint. Define
the graph
\begin{equation}\label{seed:eq:graph}
 \Phi_f(\vec x,z)\defeq
 \bigvee_{i=1}^{k}\bigl(C_i(\vec x)\land z=\tau_f(x_i,\vec x)\bigr),
\end{equation}
and $z=\tau_f(0)$ for a nullary $f$.

\begin{example}
For addition, without the general compiler:
$\Phi_+(x,y,z)\defeq(y\leq x\land z=\operatorname{ad}_x(x,y))\lor
(x<y\land z=\operatorname{ad}_y(x,y))$. In each case the chosen cap bounds
both arguments, and \cref{seed:lem:tbe-add} gives the addition equations.
Each of the two value terms is bounded by twice one input, whereas addition
has no bounding coordinate.
\end{example}

\begin{lemma}[Graphs]\label{seed:lem:tbe-graph}
$\mathsf{T}_{B\varepsilon}$ proves that each $\Phi_f$ is total and
single-valued. In the definitional extension by the functions $f^*$ with
graph $\Phi_f$:
\begin{enumerate}[label=(\alph*),nosep]
\item $f^*(\vec x)=\tau_f(b,\vec x)$ whenever $b$ bounds all $x_i$;
\item the defining equations of the description $f$ hold;
\item induction holds for quantifier-free formulas that contain finitely
many of the $f^*$.
\end{enumerate}
\end{lemma}

\begin{proof}
Exhaustiveness gives a value term in each case, and disjointness gives
uniqueness. (a) On the case $C_i$ both $b$ and $x_i$ bound all arguments, so
\cref{seed:thm:tbe-realize}(iii) applies.

(b) The initial equations follow from the initial terms. For composition
take a cap $b$ for the inputs; the inner values are bounded by $Q(b)$, and
(a) at $Q(b)$ for the outer function gives the equation from the definition
of $\tau_f$. For recursion at the counter $0$ take a cap $b$ for $0,\vec x$:
the seed is not clipped, so the value is
$\tau_g(b,\vec x)=g^*(\vec x)$. At $\Suc n$ take a cap $b$ for
$\Suc n,\vec x$: then $n\leq b$, \eqref{seed:eq:unclipped} applies, (a)
identifies the previous state with $f^*(n,\vec x)$, all arguments of the step
are bounded by $C_b$, and (a) for $h$ identifies the step value with
$h^*(n,f^*(n,\vec x),\vec x)$.

(c) Conservativity of the extension is not enough for this. By recursion on
a term with occurrences of the $f^*$ one obtains finitely many pairs
$(A_j,t_j)$ of a quantifier-free guard and a term, both in the original
language, such that the guards cover all inputs and the term equals $t_j$ on
$A_j$: at a composition take the product of the cases of the arguments and
then the cases of the outer graph. The guards need not be disjoint; the
values agree on overlaps. For a quantifier-free formula take a common
refinement for all its terms, replace the terms by their values on each
guard, and keep the connectives. The disjunction of the guarded formulas is
quantifier-free in the original language and provably equivalent to the given
formula, also under implication and negation, because the guards cover all
inputs. \Cref{seed:lem:tbe-qf} applied to it gives the induction instance. A
graph with an existential quantifier would not allow this step.
\end{proof}

\begin{theorem}[Bi-interpretation]\label{seed:thm:tied-biint}
$\mathsf{T}_{B\varepsilon}$ and $\iPRA$ are bi-interpretable. Both
interpretations use the whole domain and actual equality, the logic is
translated identically, and the identity is the isomorphism of both
composites. The theories are equiconsistent.
\end{theorem}

\begin{proof}
\emph{In $\iPRA$.} Interpret $0$, $\Suc$ as themselves, $\varepsilon$ as the
primitive recursive characteristic of equality, and each staged recursor
with body $h$ by primitive recursion with initial value
$\Suc(h(0,0,\vec p))$ and step $h(n,v,\vec p)$, after translating the body.
The equations \eqref{seed:eq:B0}, \eqref{seed:eq:BS} become defining
equations, and \eqref{seed:eq:erange}--\eqref{seed:eq:sep} and induction for
literal equations are provable.

\emph{In $\mathsf{T}_{B\varepsilon}$.} Interpret each description $f$ by
its graph $\Phi_f$, or work in the extension by the $f^*$.
\Cref{seed:lem:tbe-graph} gives the defining equations and every
quantifier-free induction instance; \cref{seed:lem:tied-macros} and
\eqref{seed:eq:sep} give injectivity of the successor and
$0\neq\Suc x$.

\emph{The composite on the side of $\iPRA$.} By recursion on a description,
with quantifier-free induction on the counter at each recursion node, the
double interpretation of a function equals the function: both satisfy the
same initial and step equations.

\emph{The composite on the side of $\mathsf{T}_{B\varepsilon}$.} By
recursion on the stage of a term. $0$ and $\Suc$ are preserved. The
interpreted characteristic of equality and $\varepsilon$ are both Boolean
and both reflect equality, so they agree by cases. For a recursor with body
$h$ the hypothesis for the earlier stage identifies the interpreted body
with $h$; the recursor and its double interpretation satisfy
\eqref{seed:eq:B0}, \eqref{seed:eq:BS} for the same body, and induction on
the equation between them identifies them at every counter and parameter.
Where the double interpretation is expressed with graph-defined functions,
\cref{seed:lem:tbe-graph}(c) provides the induction instance.

All operations and equality are preserved by the identity, so every formula
is provably equivalent to its double translation by induction on the
formula. Falsity is preserved in both directions.
\end{proof}

\subsection{The functions denoted by terms}

For $k\geq1$ let $\mathcal A_k$ be the set of primitive recursive
$f\colon\mathbb N^k\to\mathbb N$ for which there are a coordinate $i$ and a
unary primitive recursive $F$ with $f(\vec x)\leq F(x_i)$ for all $\vec x$,
and let $\mathcal A_0$ be the constants.

\begin{theorem}[Exact class of term functions]\label{seed:thm:tbe-exact}
The functions denoted in the standard model by the terms of
$\mathsf{T}_{B\varepsilon}$ with $k$ free variables are exactly the
functions in $\mathcal A_k$.
\end{theorem}

\begin{proof}
Every term denotes a primitive recursive function, and
\cref{seed:thm:tied-bound} gives the bound. Conversely let
$f(\vec x)\leq F(x_i)$. The compiler provides the cases $C_j(\vec x)$ and the
value terms $v_j(\vec x)\defeq\tau_f(x_j,\vec x)$, with value $f(\vec x)$ on
$C_j$. Let $a_j(\vec x)\defeq G(\chi_{C_j}(\vec x),v_j(\vec x))$: it is $0$
off its case and $f(\vec x)$ on it. Exactly one case holds, so each $a_j$ is
at most $F(x_i)$ and at most one is nonzero. Every unary primitive recursive
function is a term: for one variable the only case is automatic and the
value is $\tau_F(x,x)$. Put $b\defeq\tau_F(x_i,x_i)$, $s_0\defeq0$ and
$s_j\defeq\operatorname{ad}_b(s_{j-1},a_j)$ for $1\leq j\leq k$. Every
partial sum is $0$ or $f(\vec x)$, so both arguments of each bounded addition
are at most $b$, and $s_k$ denotes $f$. For $k=0$ use the numeral.
\end{proof}

The converse direction assumes that the majorant is valid in the standard
model. To prove inside $\mathsf{T}_{B\varepsilon}$ that the single term
equals the graph-defined function one needs a proof of the majorant in the
theory; \cref{seed:thm:tied-biint} needs no such hypothesis.

\begin{proposition}[Number of value terms]\label{seed:prop:tbe-sharp}
The graph \eqref{seed:eq:graph} of a $k$-ary function has $k$ value terms.
The $k$-ary sum is not, in the standard model, always equal to one of $r<k$
terms.
\end{proposition}

\begin{proof}
Each of the $r$ terms is bounded by a constant or by a function of one
input. Choose an input that controls none of them, fix the others at $0$ and
let the chosen one grow: the terms stay bounded and the sum does not.
\end{proof}

This concerns a choice among finitely many value terms. It is not a bound on
the number of clauses of a graph formula: the characteristic function of a
primitive recursive graph is bounded, hence a term by
\cref{seed:thm:tbe-exact}, so every such graph has a definition by one
equation. The value terms are what gives totality and the transfer of
induction in \cref{seed:lem:tbe-graph}.

\begin{example}[The majorant must be primitive recursive]
\label{seed:ex:tbe-majorant}
$\mathcal A_2$ is strictly smaller than the set of primitive recursive
functions that are bounded by an arbitrary function of one input. Let
$u_0,u_1,\ldots$ enumerate the unary primitive recursive descriptions and
let $g(x)\defeq1+\max_{e\leq x}u_e(x)$. This $g$ is total recursive and has
no primitive recursive majorant, since $g(x)>u_e(x)$ for $x\geq e$. Fix a
machine for $g$ and let $s(x,t)$ be its output if it halts within $t$ steps
on $x$, and $0$ otherwise; $s$ is primitive recursive. Then
$f(x,t)\defeq x+s(x,t)$ is primitive recursive, and for fixed $x$ its values
lie in $\{x,x+g(x)\}$, so it is bounded by a function of $x$. A primitive
recursive bound in $x$ would majorize $g$, and no bound in $t$ exists because
$f(x,t)\geq x$. Hence $f\notin\mathcal A_2$.
\end{example}

\subsection{Truncated subtraction in place of the equality test}

\begin{definition}\label{seed:def:tbd}
The terms are
$t\mathrel{::=}x\mid0\mid\Suc t\mid d(t,t)\mid B(n;u,v.h)$. The computation
equations are \eqref{seed:eq:B0}, \eqref{seed:eq:BS} and
\begin{equation}\label{seed:eq:monus}
 d(x,0)=x,\qquad d(x,\Suc y)=P(d(x,y)),
\end{equation}
where $P$ is the term of \cref{seed:def:tied}.
\begin{enumerate}[label=(\alph*),nosep]
\item $\mathsf{T}_{Bd}$ is ordinary intuitionistic first-order logic with
equality over these terms, with the four computation equations, coverage
$\forall x\,(x=0\lor\exists y\,x=\Suc y)$, separation $0=\Suc 0\to\bot$ and
induction for literal equations with side assumptions.
\item $\mathsf{N}_{Bd}$ is the native calculus over these terms with the
four computation equations and the rule forms of \cref{base:sec:kernel};
the contradiction is $0=\Suc 0$.
\end{enumerate}
\end{definition}

In $\mathsf{T}_{Bd}$ coverage is an axiom, in contrast to
$\mathsf{T}_{B\varepsilon}$. It is used in the proof of
\cref{seed:lem:tbd-gc}, and it does not follow from the computation
equations and induction for literal equations: \pkg{monus\_and\_neutral\_major}
refers to the coordinatewise structure on $\mathbb N\times\mathbb N$, treated
in \pkg{scope\_kernel}, which satisfies those and refutes coverage. In
$\mathsf{N}_{Bd}$ the case rule \textsc{nc} supplies this resource.

By \cref{seed:lem:tied-macros,seed:lem:tied-gate} the terms $P$, $Z$, $G$
and $B_0$ and their equations are available in both theories. In
$\mathsf{T}_{Bd}$, coverage and the values of $Z$ give
\begin{equation}\label{seed:eq:tbd-Z}
 Z(x)=0\lor Z(x)=1,\qquad Z(x)=1\leftrightarrow x=0 .
\end{equation}

\begin{lemma}\label{seed:lem:tbd-prelim}
With induction for literal equations only,
\[
 d(0,y)=0,\qquad P(d(\Suc x,y))=d(x,y),\qquad d(\Suc x,\Suc y)=d(x,y),
 \qquad d(x,x)=0,\qquad d(\Suc x,x)=1 .
\]
\end{lemma}

\begin{proof}
$d(0,y)=0$ by induction on $y$ with $P(0)=0$. For the second equation fix
$x$ and induct on $y$: the base is $P(\Suc x)=x$, and
$P(d(\Suc x,\Suc y))=P(P(d(\Suc x,y)))=P(d(x,y))=d(x,\Suc y)$. The third
equation follows from \eqref{seed:eq:monus} and the second. The last two are
inductions on $x$ with bases $d(0,0)=0$, $d(1,0)=1$ and the third equation
as step.
\end{proof}

\begin{lemma}[Gated complement]\label{seed:lem:tbd-gc}
$\mathsf{T}_{Bd}$ proves
\begin{equation}\label{seed:eq:gc}
 G\bigl(Z(d(y,x)),\,d(x,d(x,y))\bigr)=G\bigl(Z(d(y,x)),\,y\bigr),
\end{equation}
and consequently $d(y,x)=0\to d(x,d(x,y))=y$.
\end{lemma}

\begin{proof}
Fix $x$ and apply induction for literal equations to \eqref{seed:eq:gc} as
an equation in $y$. At $y=0$ the control is $Z(d(0,x))=1$, and both second
arguments are $0$, by $d(x,0)=x$ and $d(x,x)=0$.

For the step keep the induction hypothesis and put $a\defeq d(\Suc y,x)$. By
coverage $a$ is a successor or $0$. If $a$ is a successor, $Z(a)=0$ and both
sides of the goal are $0$. Let $a=0$. By \cref{seed:lem:tbd-prelim},
$d(y,x)=P(d(\Suc y,x))=0$, so the induction hypothesis, with both gates at
the control $1$, gives
\[
 d(x,t)=y,\qquad t\defeq d(x,y).
\]
If $t=0$, then $d(x,t)=x$, so $x=y$ and $a=d(\Suc y,y)=1$, contradicting
$a=0$. Hence $t=\Suc w$ for some $w$ by coverage, and
$w=P(t)=d(x,\Suc y)$ by \eqref{seed:eq:monus}. Put $v\defeq d(x,w)$. Then
$P(v)=d(x,\Suc w)=d(x,t)=y$. If $v=0$, then $y=0$, so
$t=d(x,0)=x=\Suc w$ and $v=d(\Suc w,w)=1$, a contradiction. Hence
$v=\Suc z$ by coverage, $z=P(v)=y$ and $v=\Suc y$. With $w=d(x,\Suc y)$
this reads
\[
 d(x,d(x,\Suc y))=\Suc y,
\]
which is the goal at the control $Z(a)=1$. The contradictions use
$0\neq1$ only. The only induction formula is the literal equation
\eqref{seed:eq:gc}. If $d(y,x)=0$, both gates have the control $1$.
\end{proof}

\begin{theorem}[Reflection of equality]\label{seed:thm:tbd-reflect}
Let $\varepsilon(x,y)\defeq G\bigl(Z(d(x,y)),Z(d(y,x))\bigr)$.
$\mathsf{T}_{Bd}$ proves
\[
 d(x,y)=0\land d(y,x)=0\to x=y,\qquad
 \varepsilon(x,y)=0\lor\varepsilon(x,y)=1,\qquad
 \varepsilon(x,y)=1\leftrightarrow x=y .
\]
\end{theorem}

\begin{proof}
From $d(y,x)=0$, \cref{seed:lem:tbd-gc} gives $d(x,d(x,y))=y$; substituting
$d(x,y)=0$ and $d(x,0)=x$ gives $x=y$. Each input of the gate is $0$ or $1$
by \eqref{seed:eq:tbd-Z}, and the gate equations give the range and
$\varepsilon(x,y)=1\leftrightarrow d(x,y)=0\land d(y,x)=0$: if the first
input is $0$ the output is $0$, and otherwise it is the second input. The
first statement gives $x=y$ from the right side; conversely $x=y$ makes both
differences $0$ by $d(x,x)=0$.
\end{proof}

\begin{corollary}\label{seed:cor:tbd-biint}
$\mathsf{T}_{Bd}$ is bi-interpretable with $\iPRA$, with whole domains,
actual equality, the logic translated identically and the identity as
isomorphism of both composites.
\end{corollary}

\begin{proof}
By \cref{seed:thm:tbd-reflect} the term $\varepsilon$ satisfies
\eqref{seed:eq:erange} and \eqref{seed:eq:ereflect}, so every derivation of
$\mathsf{T}_{B\varepsilon}$ becomes a derivation of $\mathsf{T}_{Bd}$ when
this term is substituted for the symbol; in particular
\cref{seed:thm:tbe-realize,seed:lem:tbe-graph} hold. In the other direction
$\iPRA$ interprets the recursors as in \cref{seed:thm:tied-biint} and $d$ as
truncated subtraction; it proves \eqref{seed:eq:monus}, coverage, separation
and induction for literal equations.

For the composite on the side of $\mathsf{T}_{Bd}$, first the recursors
$C$ and $W$, which do not contain $d$, are identified with their double
interpretations by induction on the equation between them, and hence $P$ is
preserved. The double interpretation of $d$ is the graph of the primitive
recursive truncated subtraction; it satisfies the two equations
\eqref{seed:eq:monus} with the preserved $P$, and induction on the second
argument identifies it with $d$. Likewise the term $\widehat d$ of
\cref{seed:lem:tbe-diff}, with $\varepsilon$ expanded, satisfies
$\widehat d(x,0)=x$ and $\widehat d(x,\Suc y)=P(\widehat d(x,y))$, so
$\widehat d(x,y)=d(x,y)$ by induction on $y$ for this literal equation of
the original language. The remaining recursors are identified as in
\cref{seed:thm:tied-biint}, and the composite on the side of $\iPRA$ is
unchanged.
\end{proof}

\subsection{The native calculus with truncated subtraction}

The proof of \cref{seed:lem:tbd-gc} splits on $t=0$ or $t=\Suc w$ while an
induction hypothesis that mentions $t$ is in the context, and
\cref{seed:cor:tied-noterms} excludes a selector term. $\mathsf{N}_{Bd}$
derives the split by a selector whose zero branch is Boolean.

\begin{lemma}\label{seed:lem:nbd-basic}
$\mathsf{N}_{Bd}$ derives the equality laws, \textsc{nc} and \textsc{ei} of
\cref{seed:lem:succ-basic}, the equations of
\cref{seed:lem:tied-macros,seed:lem:tied-gate,seed:lem:tbd-prelim}, every
equation from $0=\Suc 0$, and the functional range rule: if
$\Gamma\proves K(0)$ and $\Gamma\proves K(1)$ then
$\Gamma\proves K(Z(t))$, for every formula template $K$ and term $t$.
\end{lemma}

\begin{proof}
Reflexivity follows from the computation $C(\Suc 0,p)=p$ as in
\cref{seed:lem:succ-basic}, and the other equality laws, \textsc{nc} and
\textsc{ei} are derived as there; \cref{seed:lem:succ-leibniz} holds as
well. The equations of the three lemmas use computation, equality and
\textsc{ei} only. Under $0=\Suc 0$, transport between $C(0,a)=\Suc a$ and
$C(1,a)=a$ gives $\Suc a=a$; then \textsc{ei} on the equation $a=0$ gives
$a=0$ for every $a$, hence every equation. The functional range rule is
\textsc{nc} on the template $w\mapsto K(Z(w))$.
\end{proof}

By recursion on a native formula, with \rulename{DI} by cases at each
$\D{x}{y}$ and the representation of $\Q{y}B$ by $\D{v}{y}B$ with $v$ unused,
every native formula follows from $0=\Suc 0$.

A term $b$ is called \emph{two-valued} in $\Gamma$ if the functional range
rule holds for it: $\Gamma\proves K(0)$ and $\Gamma\proves K(1)$ imply
$\Gamma\proves K(b)$ for every template $K$. Terms $Z(t)$ are two-valued.

\begin{lemma}[A selector with a Boolean zero branch]\label{seed:lem:nbd-sel}
Let $\operatorname{inc}(a,b)\defeq C(Z(b),a)$ and
\[
 \eta(y,b,u)\defeq\operatorname{inc}\bigl(G(y,u),\,G(Z(y),b)\bigr).
\]
Then $\eta(\Suc y,b,u)=u$ for all terms, and $\eta(0,b,u)=b$ for every
two-valued $b$.
\end{lemma}

\begin{proof}
$\eta(\Suc y,b,u)=\operatorname{inc}(u,G(0,b))=\operatorname{inc}(u,0)
=C(1,u)=u$. $\eta(0,b,u)=\operatorname{inc}(0,G(1,b))=C(Z(b),0)$, which is
$C(1,0)=0$ at $b=0$ and $C(0,0)=1$ at $b=1$; the functional range rule for
the template $C(Z(w),0)=w$ gives $\eta(0,b,u)=b$.
\end{proof}

The zero branch of $\eta$ is two-valued and the successor branch is
arbitrary, so $\eta$ does not contradict \cref{seed:thm:tied-bound}.

\begin{lemma}[Contextual cases]\label{seed:lem:nbd-cases}
In $\mathsf{N}_{Bd}$, let $s=t$ be a literal equation.
\begin{enumerate}[label=(\alph*),nosep]
\item If $b$ is two-valued in $\Gamma$, $\Gamma,b=0\proves s=t$ and
$\Gamma,b=1\proves s=t$, then $\Gamma\proves s=t$.
\item If $z$ is fresh, $\Gamma,r=0\proves s=t$ and
$\Gamma,r=\Suc z\proves s=t$, then $\Gamma\proves s=t$.
\end{enumerate}
The context and the equation may mention $b$ and $r$.
\end{lemma}

\begin{proof}
(a) Let $K(x)\defeq\D{a}{y}\bigl(\eta(y,x,s)=\eta(y,b,t)\bigr)$ with $a$
unused. For $c\in\{0,1\}$ derive $K(c)$ by \rulename{DI} by cases: the local
assumption $\eta(0,c,s)=\eta(0,b,t)$ is $c=b$ by \cref{seed:lem:nbd-sel},
the given branch gives $s=t$, and this is the matrix at $\Suc z$. The range
rule for $b$ gives $K(b)$. Rule \rulename{DE} opens its closure, whose zero
instance $\eta(0,b,s)=\eta(0,b,t)$ holds because both sides equal $b$, and
\rulename{QE} at $1$ gives $s=t$.

(b) Rule \textsc{nc} on fresh variables in literal equations gives
\[
 G(Z(r),r)=0,\qquad G(Z(Z(r)),r)=G\bigl(Z(Z(r)),\Suc(P(r))\bigr):
\]
for the first, the zero case is $G(1,0)=0$ and a successor case has the
control $0$; for the second, the zero case has the control $0$, and a
successor case has the control $1$ and $P(\Suc x)=x$. Apply (a) to
$b\defeq Z(r)$. Under $Z(r)=1$ the first identity gives $r=0$. Under
$Z(r)=0$ the second gives $r=\Suc(P(r))$; substitute $P(r)$ for $z$ in the
second given branch.
\end{proof}

\begin{theorem}\label{seed:thm:nbd}
$\mathsf{N}_{Bd}$ derives \eqref{seed:eq:gc}, the three statements of
\cref{seed:thm:tbd-reflect} in the forms
\[
 d(x,y)=0,\ d(y,x)=0\proves x=y,\qquad
 s=t\proves\varepsilon(s,t)=1,\qquad\varepsilon(s,t)=1\proves s=t,
\]
with $\varepsilon(s,t)$ two-valued, and the introduction and elimination
rules of intuitionistic implication, conjunction, disjunction, both
quantifiers and falsity for the formulas defined in the proof. It is
bi-interpretable with $\mathsf{T}_{Bd}$ and with $\iPRA$, with whole domains,
actual equality, identity composites and both formula round trips.
\end{theorem}

\begin{proof}
Every case split in the proof of \cref{seed:lem:tbd-gc} has an equation as
its conclusion, so it is an instance of \cref{seed:lem:nbd-cases}(b); the
branches that end in $0=1$ are closed by \cref{seed:lem:nbd-basic}. The
induction is \textsc{ei} for the same literal equation. The term
$\varepsilon(s,t)$ is two-valued by the range rules of its two inputs and
the gate equations.

For a literal equation $E$, namely $s=t$, define the gate
$\mathcal G_v(E)\defeq\bigl(G(v,\varepsilon(s,t))=G(v,1)\bigr)$; it is
derivable at $v=0$ and interderivable with $E$ at $v=\Suc w$. For literal
equations $E$ ($s=t$) and $F$ ($r=w$) define the formula selector
$\mathcal M_v(E,F)\defeq\bigl(\eta(v,\varepsilon(s,t),\varepsilon(r,w))=1
\bigr)$; by \cref{seed:lem:nbd-sel}, $\mathcal M_0(E,F)$ is interderivable
with $E$ and $\mathcal M_{\Suc z}(E,F)$ with $F$. Extend the gate through
$\D{x}{y}$ by $\mathcal G_v(\D{x}{y}A)\defeq\D{x}{y}\mathcal G_vA$, replace
each $\Q{y}A$ by $\D{a}{y}A$ with $a$ unused, align binder depths by the
padding $A\approx\D{a}{y}\mathcal G_yA$, and extend the selector through
matched binders by recursion on the formulas; the endpoint properties are
preserved by the binder congruence of \cref{seed:lem:succ-leibniz}. With
$p(z)\defeq B_0(z;u,v.Z(v))$, so that $p(1)=1$ and $p(2)=0$, and with unused
coordinates fresh, put
\begin{align*}
 A\Rightarrow_*B&\defeq\D{a}{z}\mathcal M_z(A,B), &
 A\land_*B&\defeq\D{a}{z}\mathcal G_z\mathcal M_{p(z)}(A,B),\\
 \exists_*x\,A(x)&\defeq\D{x}{y}\mathcal G_yA(x), &
 A\lor_*B&\defeq\exists_*x\,\mathcal M_x(A,B),\\
 \forall_*x\,A(x)&\defeq A(0)\land_*\D{a}{z}\mathcal G_zA(z), &
 \bot_*&\defeq(0=\Suc 0).
\end{align*}
Implication introduction is \rulename{DI} by cases, whose local assumption
is the zero instance of the selector; implication elimination opens the
closure and specializes at $1$. The derivable zero instance of the gate
gives the introductions of conjunction and existence, and specialization at
$1$ and $2$ their eliminations. Once implication is available, \textsc{nc}
on an implication template gives numerical cases in a context for arbitrary
formulas, hence disjunction elimination and universal elimination. These are
the constructions of \cref{base:thm:bridge} with the atomic gate and
selector above in place of the numerical selector $\sel$, which does not
exist here.

Every rule of $\mathsf{N}_{Bd}$ is sound for $\mathsf{T}_{Bd}$ under the
readings of \cref{base:sec:reading}: \rulename{DI} by cases uses coverage
and \rulename{QI} uses induction for literal equations. Conversely
$\mathsf{N}_{Bd}$ derives the translation of every theorem of
$\mathsf{T}_{Bd}$ by the rules just established. The round trips are
\[
 \Q{y}B\approx B(0)\Rightarrow_*\forall_*y\,B(y),\qquad
 \D{x}{y}B\approx
 \exists_*x\,\bigl(B(x,0)\Rightarrow_*\forall_*y\,B(x,y)\bigr):
\]
forward, \rulename{DE} opens the witness, and under the zero instance
\rulename{QE} gives every instance; backward, eliminate the existential, keep
the zero instance of the matrix, and use the implication and the universal
at $\Suc z$ for \rulename{DI} by cases. Induction on formulas gives both
formula composites, and the numerical composites are identities. The
bi-interpretation with $\iPRA$ is \cref{seed:cor:tbd-biint}.
\end{proof}

%======================================================================
\section{The self-seeded recursor}
\label{seed:sec:self}

This section draws on \pkg{selfseed\_scalar}, \pkg{scope\_and\_coverage} and
\pkg{coverage\_and\_kernel}.

\subsection{The theories}

\begin{definition}\label{seed:def:tf}
The terms are
$t\mathrel{::=}x\mid0\mid\Suc t\mid\varepsilon(t,t)\mid F(n;v.h)$, where
$F(n;v.h)$ binds $v$ in $h$. The computation equations are
\begin{equation}\label{seed:eq:F}
 F(0;v.h)=\Suc(h\subst{0}{v}),\qquad
 F(\Suc n;v.h)=h\subst{F(n;v.h)}{v}.
\end{equation}
$\mathsf{T}_F^{-}$ is ordinary intuitionistic first-order logic with
equality over these terms with \eqref{seed:eq:F}, the axioms
\eqref{seed:eq:erange}, \eqref{seed:eq:ereflect}, \eqref{seed:eq:sep} for
$\varepsilon$ and separation, and induction for literal equations with side
assumptions. $\mathsf{T}_F$ is $\mathsf{T}_F^{-}$ with the additional axiom
$\Suc x=\Suc y\to x=y$. Coverage is not an axiom of either theory.
\end{definition}

The update sees neither an initializer nor the old counter. Compared with
$\mathsf{T}_{B\varepsilon}$, which derives injectivity from the counter
argument (\cref{seed:lem:tied-macros}), $\mathsf{T}_F$ postulates it. As in
\cref{seed:sec:tied}, $\Suc$ can be removed from the signature by
$\Suc a\defeq F(0;v.a)$. Put
\[
 C(n,a)\defeq F(n;v.a),\qquad Z(n)\defeq C(n,0),
\]
so that $C(0,a)=\Suc a$, $C(\Suc n,a)=a$, $Z(0)=1$, $Z(\Suc n)=0$.

\begin{lemma}\label{seed:lem:tf-basic}
$\mathsf{T}_F^{-}$ proves $\Suc n\neq0$. The terms
\[
 \neg_Bb\defeq\varepsilon(b,0),\qquad
 b\land_Bc\defeq\varepsilon\bigl(C(b,\varepsilon(c,0)),0\bigr),\qquad
 b\lor_Bc\defeq\neg_B(\neg_Bb\land_B\neg_Bc)
\]
have the truth tables of negation, conjunction and disjunction on the values
$0,1$. Every quantifier-free formula $A$ has a term $\chi_A$ with
$\chi_A=0\lor\chi_A=1$ and $\chi_A=1\leftrightarrow A$, and
$\mathsf{T}_F^{-}$ admits induction with side assumptions for
quantifier-free formulas.
\end{lemma}

\begin{proof}
If $\Suc n=0$, then $0=Z(\Suc n)=Z(0)=1$. Reflection and the range axiom
give $\varepsilon(0,0)=1$ and $\varepsilon(1,0)=0$. For the conjunction,
$C(0,\varepsilon(c,0))=\Suc(\varepsilon(c,0))\neq0$, so the value at $b=0$
is $0$; and $C(1,\varepsilon(c,0))=\varepsilon(c,0)$, so the value at $b=1$
is $\neg_B\neg_Bc=c$. The only non-Boolean intermediate value is a
successor, and its distinctness from $0$ does not cancel a successor. Define
$\chi_{s=t}\defeq\varepsilon(s,t)$, $\chi_\bot\defeq0$, and the connectives
by the Boolean terms, implication as $\neg_B\chi_A\lor_B\chi_B$; the Boolean
range comes from the range axiom and the truth tables. Induction for $A(n)$
is induction for the literal equation $\chi_{A(n)}=1$.
\end{proof}

\begin{definition}\label{seed:def:finite-case}
A \emph{finite-case graph} is a formula
$\bigvee_i(A_i(\vec x)\land y=t_i(\vec x))$ with quantifier-free guards $A_i$
and terms $t_i$ of the language, such that the guards are provably
exhaustive and the values provably agree where guards overlap.
\end{definition}

\begin{lemma}[Finite-case elimination]\label{seed:lem:tf-cases}
A finite-case graph is provably total and single-valued. In the definitional
extension by the functions of finitely many finite-case graphs, every
quantifier-free formula is provably equivalent to a quantifier-free formula
of the original language, and induction for quantifier-free formulas of the
extension is admissible. Compositions of finite-case functions are
finite-case.
\end{lemma}

\begin{proof}
As in \cref{seed:lem:tbe-graph}(c): expand each term into finitely many
guarded terms of the original language by refining the cases of its
arguments and then the cases of the outer graph, take a common refinement
for a formula, and apply \cref{seed:lem:tf-basic} to the resulting
disjunction.
\end{proof}

The lemma concerns function terms in formulas. It does not allow a new body
$h$ of $F$ that contains a function defined by a graph; bodies are always
terms of the original language.

\subsection{Without injectivity of the successor}

All statements of this subsection are theorems of $\mathsf{T}_F^{-}$.

\begin{lemma}[No fixed point]\label{seed:lem:tf-nofix}
$\Suc a\neq a$ for every $a$.
\end{lemma}

\begin{proof}
The term $h(v)\defeq\varepsilon(v,0)$ has no fixed point: $h(v)=v$ and the
range of $h$ give $v=0$ or $v=1$, contradicting $h(0)=1$ and $h(1)=0$. If
$\Suc a=a$, then $F(a;v.h)=F(\Suc a;v.h)=h(F(a;v.h))$.
\end{proof}

\begin{theorem}[No cycle of a fixed length]\label{seed:thm:tf-acyclic}
For each fixed integer $k\geq1$, the equations \eqref{seed:eq:F} and the
axioms \eqref{seed:eq:erange}, \eqref{seed:eq:ereflect}, \eqref{seed:eq:sep}
prove $\forall a\;\Suc^ka\neq a$, without any instance of induction.
Consequently distinct numerals are provably distinct, and
$\Suc^ia\neq\Suc^ja$ for fixed $i<j$.
\end{theorem}

\begin{proof}
The proofs are constructed by recursion on $k$ outside the theory. Assume
the statement for all lengths below $k$. Then $\bar0,\ldots,\bar k$ are
pairwise distinct: $\bar0$ differs from a positive numeral by
\cref{seed:lem:tf-basic}, and $\bar i=\bar j$ with $1\leq i<j\leq k$ would
be $\Suc^{j-i}(\bar i)=\bar i$ with $1\leq j-i<k$. For $k=1$ only separation
is used.

For a term $b$ with $b=0\lor b=1$ let
$\operatorname{Inc}(b,a)\defeq C(\varepsilon(b,0),a)$, so that
$\operatorname{Inc}(0,a)=a$ and $\operatorname{Inc}(1,a)=\Suc a$. Define, as
a written finite expansion,
\[
 T_0(v)\defeq0,\qquad
 T_{j+1}(v)\defeq\operatorname{Inc}\bigl(\varepsilon(v,\bar j),\cdot\bigr)^{j+1}
 \bigl(T_j(v)\bigr)\quad(0\leq j<k),\qquad c_k(v)\defeq T_k(v),
\]
where the exponent denotes $j+1$ written compositions. At $v=\bar i$ with
$i<k$ the test with index $i$ is $1$ and the others are $0$, so $i+1$
increments are performed; at $v=\bar k$ all tests are $0$. Hence
\[
 c_k(\bar i)=\overline{i+1}\ (0\leq i<k),\qquad c_k(\bar k)=0,\qquad
 c_k(v)=\bar0\lor\cdots\lor c_k(v)=\bar k,
\]
the last by the finitely many cases of \eqref{seed:eq:erange}: if some test
is $1$, reflection identifies $v$ with a numeral and distinctness makes the
other tests $0$.

If $c_k^{\,k}(w)=w$ for the $k$-fold composition, then $w$ is one of
$\bar0,\ldots,\bar k$ by the range, since $k\geq1$. On these $k+1$ distinct
values $c_k$ is one cycle, and $k$ applications of it end at a different
numeral. So $c_k^{\,k}(w)\neq w$ for every $w$.

Now let $\Suc^ka=a$ and $w\defeq F(a;v.c_k)$. By $k$ uses of the successor
equation and congruence,
$w=F(\Suc^ka;v.c_k)=c_k^{\,k}(F(a;v.c_k))=c_k^{\,k}(w)$, a contradiction.
The last statement applies the theorem to $\Suc^ia$ with the length $j-i$.
\end{proof}

The theorem excludes a cycle along one successor chain. It does not derive
$a=b$ from $\Suc a=\Suc b$: two predecessor chains may merge without a
cycle. The term $c_k$ depends on the fixed $k$, and there is no term for
cycles of a length given by a variable.

For the next results define, with a parameter $a$,
\[
 g_a(v)\defeq\varepsilon(v,0)\lor_B\varepsilon(\Suc v,a),\qquad
 \theta_a(v)\defeq C(g_a(v),v),\qquad
 U(j,a)\defeq F(j;v.\theta_a(v)),\qquad p_*(a)\defeq U(a,a).
\]
The update $\theta_a$ holds the state if it is $0$ or if its successor is
$a$, and increments it otherwise. Since $\theta_a(0)=0$, the orbit starts at
$U(0,a)=1$, and $U(\Suc j,a)=\theta_a(U(j,a))$.

\begin{lemma}[Protected orbit]\label{seed:lem:tf-orbit}
$U(j,a)=\Suc j\lor\Suc(U(j,a))=a$, and $\Suc(p_*(\Suc^2m))=\Suc^2m$.
\end{lemma}

\begin{proof}
The first statement is a quantifier-free induction on $j$ with $a$ fixed.
The base is $U(0,a)=\Suc 0$. In the step put $v\defeq U(j,a)$ and decide
$\Suc v=a$. If it holds, $\theta_a$ fixes $v$ and the second alternative
persists. Otherwise the first alternative of the hypothesis gives
$v=\Suc j\neq0$, the guard is false and $U(\Suc j,a)=\Suc v=\Suc\Suc j$.

Specialize the proved statement to $j=m$, $a=\Suc^2m$. Both alternatives
give $\Suc(U(m,\Suc^2m))=\Suc^2m$, so the state is fixed by $\theta_a$ and
two further computation steps give $U(\Suc^2m,\Suc^2m)=U(m,\Suc^2m)$. No
successor is cancelled.
\end{proof}

\begin{theorem}[Coverage and a right predecessor]\label{seed:thm:tf-coverage}
Let $P$ be the function of the finite-case graph
\[
 (n=0\land p=0)\lor(n=1\land p=0)\lor(n\neq0\land n\neq1\land p=p_*(n)).
\]
$\mathsf{T}_F^{-}$ proves, by \cref{seed:lem:tf-cases},
\[
 P(0)=P(1)=0,\qquad n=0\lor n=\Suc(P(n)),\qquad\Suc(P(\Suc n))=\Suc n,
\]
and hence coverage $n=0\lor\exists m\,n=\Suc m$.
\end{theorem}

\begin{proof}
The graph is finite-case because equality is decidable and $0\neq1$. Apply
quantifier-free induction to
\[
 n=0\ \lor\ n=1\ \lor\ \Suc(p_*(n))=n .
\]
The base holds. In the step, $n=0$ gives $\Suc n=1$. If $n=1$, then
$\Suc n=\Suc^20$ and \cref{seed:lem:tf-orbit} with $m=0$ gives
$\Suc(p_*(\Suc n))=\Suc n$. If $\Suc(p_*(n))=n$, then
$\Suc n=\Suc^2(p_*(n))$ and \cref{seed:lem:tf-orbit} with $m=p_*(n)$ gives
$\Suc(p_*(\Suc n))=\Suc n$. This proves the displayed disjunction for every
$n$. To pass to $P$, decide $n=0$ and $n=1$: if $n=0$ the first alternative
of $n=0\lor n=\Suc(P(n))$ holds; if $n=1$ then $n=\Suc 0=\Suc(P(1))$; and if
$n\neq0,1$ the third alternative of the disjunction holds and
$P(n)=p_*(n)$, so $n=\Suc(P(n))$. Applying $n=0\lor n=\Suc(P(n))$ to
$\Suc n$ and excluding its zero alternative gives
$\Suc(P(\Suc n))=\Suc n$.
\end{proof}

The proof holds for all elements of a model, standard or not. The result is
a right inverse: it chooses one predecessor of each nonzero element and does
not prove that $P(\Suc a)=a$.

\begin{proposition}[Standard numerals have unique predecessors]
\label{seed:prop:tf-rigid}
For every fixed integer $k\geq1$, $\mathsf{T}_F^{-}$ proves
\begin{equation}\label{seed:eq:rigid}
 \Suc x=\bar k\leftrightarrow x=\overline{k-1}.
\end{equation}
\end{proposition}

\begin{proof}
By recursion on $k$ outside the theory. The data carried are the pairwise
distinctness of $\bar0,\ldots,\overline{k+1}$ and \eqref{seed:eq:rigid} for
all smaller positive indices. Initially $0,1,2$ are distinct by
\cref{seed:lem:tf-basic,seed:lem:tf-nofix}.

Let $h_k$ be a term with the values
\[
 h_k(v)=\begin{cases}
 \overline{i-1},& v=\bar i,\ 1\leq i\leq k-1,\\
 \overline{k-1},& v=\overline{k+1},\\
 \bar k,&\text{otherwise.}
 \end{cases}
\]
Such a term is obtained as follows. For $r=1,\ldots,k$ let $b_r(v)$ be the
Boolean negation of the disjunction of $\varepsilon(v,\bar i)$ for
$1\leq i\leq\min(r,k-1)$, together with
$\varepsilon(v,\overline{k+1})$ if $r=k$; it is the characteristic term of
the condition that the displayed value is at least $r$. Put $t_0(v)\defeq0$,
$t_r(v)\defeq C(\neg_Bb_r(v),t_{r-1}(v))$ and $h_k\defeq t_k$: a true
condition contributes one successor. For $k=1$ the values of $h_1$ are those
of $\neg_B\varepsilon(v,2)$.

Let $A_k(n)\defeq F(n;v.h_k)$ and let $q_0\defeq\overline{k+1}$,
$q_j\defeq\overline{k-j}$ for $1\leq j\leq k$; these numerals are distinct.
Since $h_k(0)=\bar k$, $A_k(0)=q_0$, and computation gives
$A_k(\bar j)=q_j$ for $0\leq j\leq k$. The term $h_k$ never has the value
$q_0$, and for $1\leq j\leq k$,
\[
 h_k(v)=q_j\leftrightarrow v=q_{j-1}.
\]
Quantifier-free induction gives $A_k(n)=q_0\leftrightarrow n=0$: at a
successor both sides are false. For $1\leq j<k$, recursion on $j$ gives
$A_k(n)=q_j\leftrightarrow n=\bar j$: by coverage
(\cref{seed:thm:tf-coverage}) $n=0$ or $n=\Suc m$; at $0$ both sides are
false; at $\Suc m$ the left side is equivalent to $A_k(m)=q_{j-1}$, hence to
$m=\overline{j-1}$, and so is the right side by \eqref{seed:eq:rigid} for
$j<k$.

Now let $\Suc x=\bar k$. Then $h_k(A_k(x))=A_k(\bar k)=q_k$, so
$A_k(x)=q_{k-1}$ and $x=\overline{k-1}$. The converse is immediate. For the
distinctness of $\overline{k+2}$: an equality $\overline{k+2}=\bar j$ with
$1\leq j\leq k$ reduces by \eqref{seed:eq:rigid} for $j$ to
$\overline{k+1}=\overline{j-1}$; equality with $0$ is excluded by
\cref{seed:lem:tf-basic} and equality with $\overline{k+1}$ by
\cref{seed:lem:tf-nofix}.
\end{proof}

This is a family of derivations, one for each $k$, and not one theorem that
quantifies over $k$. It does not exclude a collision
$\Suc x=\Suc y$, $x\neq y$, over a nonstandard target.

\begin{lemma}[An equivalent form of injectivity]\label{seed:lem:tf-retract}
In the extension by $P$ put $T(n)\defeq P(\Suc n)$. Then
\[
 \Suc(T(n))=\Suc n,\qquad T(T(n))=T(n),\qquad T(0)=0,
\]
and over $\mathsf{T}_F^{-}$ injectivity of the successor is equivalent to
$\forall n\,(T(n)=n\to T(\Suc n)=\Suc n)$.
\end{lemma}

\begin{proof}
The first equation is \cref{seed:thm:tf-coverage}; applying $P$ to it gives
the second; the third is $P(1)=0$. If the fixed points of $T$ are closed
under successor, quantifier-free induction on $T(n)=n$ gives $T(n)=n$ for
all $n$, and from $\Suc a=\Suc b$ one gets $a=T(a)=T(b)=b$. Conversely
injectivity applied to $\Suc(T(n))=\Suc n$ gives $T(n)=n$ for all $n$.
\end{proof}

\begin{remark}\label{seed:rem:tf-open}
\pkg{scope\_and\_coverage} and \pkg{coverage\_and\_kernel} contain the
protected orbit with the same guard, under the names $V$, $W$ and $U$,
$p_*$; the first derives $n=0\lor n=\Suc(P(n))$ by induction on that formula
in the extension by $P$, the second by the induction of the proof above.
\Cref{seed:thm:tf-acyclic,seed:lem:tf-retract} are from the first package
and \cref{seed:lem:tf-nofix,seed:prop:tf-rigid} from the second. Neither
derives injectivity. The first package records why a comparison of the
orbits with targets $\Suc n$ and $\Suc^2n$ does not close the gap: for
$v\neq0$ the term $\Suc(\theta_{\Suc n}(v))$ stops when
$\Suc v=\Suc n$ and $\theta_{\Suc^2n}(\Suc v)$ stops when
$\Suc^2v=\Suc^2n$, so the intertwining of the two orbits needs
$\Suc^2v=\Suc^2n\to\Suc v=\Suc n$ at the visited states, which is an instance
of the missing cancellation.
\end{remark}

\subsection{With injectivity: predecessor and order}

From here on the theory is $\mathsf{T}_F$.

\begin{theorem}[Predecessor]\label{seed:thm:tf-pred}
$\mathsf{T}_F$ proves $p_*(\Suc^2m)=\Suc m$, and for the finite-case
function $P$ of \cref{seed:thm:tf-coverage}
\[
 P(0)=0,\qquad P(\Suc n)=n,\qquad n=0\lor n=\Suc(P(n)).
\]
\end{theorem}

\begin{proof}
Injectivity applied to \cref{seed:lem:tf-orbit} gives the first equation. By
coverage $n=0$ or $n=\Suc m$. In the first case $P(\Suc n)=P(1)=0$. In the
second $\Suc^2m$ differs from $0$ and, by injectivity and
$\Suc m\neq0$, from $1$, so $P(\Suc^2m)=p_*(\Suc^2m)=\Suc m$.
\end{proof}

Write $w(a)\defeq p_*(a)$. The constructions below use this term only
through
\begin{equation}\label{seed:eq:W}
 w(\Suc^3m)=\Suc^2m ,
\end{equation}
% ed.: selfseed_scalar uses its own term W (see the remark below) in the
% constructions that follow; they depend on W only through this equation,
% which p_* also satisfies, so the term of coverage_and_kernel is used
% throughout and the predecessor is proved once.
that is, $w$ removes one successor from every value that is at least $3$,
and $w(w(\cdot))$ removes two from every value that is at least $4$.

\begin{remark}\label{seed:rem:tf-W}
This order of proof, coverage first and then the predecessor with
corrections at $0$ and $1$, is that of \pkg{coverage\_and\_kernel}.
\pkg{selfseed\_scalar} uses the orbit without the guard at zero,
$V(j,a)\defeq F(j;v.C(\varepsilon(\Suc v,a),v))$ and $W(a)\defeq V(a,a)$,
with the invariant
$a\neq1\to(V(j,a)=\Suc^2j\lor\Suc(V(j,a))=a)$, which gives
\eqref{seed:eq:W} for $W$ and the values $W(0)=2$, $W(1)=2$, $W(2)=4$. Its
predecessor graph has the exceptional inputs $0,1,2$, and predecessor and
coverage are proved together by quantifier-free induction on
$(n=0\lor n=\Suc(P(n)))\land P(\Suc n)=n$. Both proofs cancel a successor at
the step from $\Suc(V(m,a))=a$ to the value of $V(m,a)$, which is where
$\mathsf{T}_F$ uses its injectivity axiom.
\end{remark}

For a parameter $a$ define the capped counter and the comparison term
\[
 \lambda_a(v)\defeq C(\varepsilon(v,\Suc^2a),v),\qquad
 \operatorname{Cap}(n,a)\defeq F(n;v.\lambda_a(v)),\qquad
 L(a,n)\defeq\varepsilon(\operatorname{Cap}(n,a),\Suc^2a),
\]
and write $a\leq n$ for $L(a,n)=1$. The update $\lambda_a$ holds the value
$\Suc^2a$ and increments every other value; since $\Suc^2a\neq0$,
$\operatorname{Cap}(0,a)=2$.

\begin{lemma}[Capped counting and order]\label{seed:lem:tf-order}
$\mathsf{T}_F$ proves
\begin{gather*}
 \operatorname{Cap}(n,a)=\Suc^2n\lor\operatorname{Cap}(n,a)=\Suc^2a,\qquad
 \operatorname{Cap}(a,a)=\Suc^2a,\qquad
 \operatorname{Cap}(\Suc n,\Suc a)=\Suc(\operatorname{Cap}(n,a)),\\
 L(0,n)=1,\qquad L(\Suc a,0)=0,\qquad L(\Suc a,\Suc n)=L(a,n),\\
 a\leq\Suc b\leftrightarrow(a\leq b\lor a=\Suc b),
\end{gather*}
and that $\leq$ is a discrete linear order with least element $0$ which the
successor preserves and reflects. The finite-case function $\min(n,a)$, with
value $n$ if $n\leq a$ and $a$ otherwise, satisfies
$\operatorname{Cap}(n,a)=\Suc^2\min(n,a)$.
\end{lemma}

\begin{proof}
The first statement is a quantifier-free induction on $n$ with $a$ fixed:
the base is $\operatorname{Cap}(0,a)=\Suc^20$; a state $\Suc^2a$ is held;
and a state $\Suc^2n\neq\Suc^2a$ is incremented to $\Suc^2\Suc n$. At $n=a$
both alternatives give the second statement. In particular
$\operatorname{Cap}(n,0)=2$.

Coverage gives $C(c,\Suc v)=\Suc(C(c,v))$, and reflection with injectivity
gives $\varepsilon(\Suc v,\Suc^3a)=\varepsilon(v,\Suc^2a)$; hence
$\lambda_{\Suc a}(\Suc v)=\Suc(\lambda_a(v))$. Also
$\operatorname{Cap}(1,\Suc a)=\lambda_{\Suc a}(2)=3$, because
$\Suc^3a\neq2$. Induction on $n$ for the literal equation, with $a$ fixed,
gives the third statement: the base is $3=\Suc 2$ and the step applies
$\lambda_{\Suc a}$ and the last identity.

$L(0,n)=1$ by $\operatorname{Cap}(n,0)=2$; $L(\Suc a,0)=\varepsilon(2,
\Suc^3a)=0$; and
$L(\Suc a,\Suc n)=\varepsilon(\Suc\operatorname{Cap}(n,a),\Suc\Suc^2a)
=L(a,n)$.

For the equivalence, let $\operatorname{Cap}(\Suc b,a)=\Suc^2a$. If
$\operatorname{Cap}(b,a)=\Suc^2a$, then $a\leq b$. Otherwise
$\operatorname{Cap}(b,a)=\Suc^2b$ by the first statement, the next value is
$\Suc^3b$, and $\Suc^3b=\Suc^2a$ gives $a=\Suc b$. Conversely a reached cap
persists, and $\operatorname{Cap}(\Suc b,\Suc b)=\Suc^2\Suc b$.

\emph{Linearity.} Fix $a$ and induct on $n$ in $a\leq n\lor n\leq a$. The
base is $0\leq a$. In the step, $a\leq n$ persists. If $n\leq a$, split
$a=0$ or $a=\Suc b$ by coverage. If $a=0$, then
$\operatorname{Cap}(0,n)=2=\Suc^2n$ gives $n=0$ and $a\leq\Suc n$. If
$a=\Suc b$, the equivalence gives $n\leq b$, and then
$\Suc n\leq\Suc b=a$, or $n=a$, and then $a\leq\Suc n$.

\emph{Irreflexivity and descent.} $\Suc n\not\leq n$ by induction on the
literal equation $L(\Suc n,n)=0$. $\Suc n\leq a\to n\leq a$: by coverage
$a=0$ contradicts $L(\Suc n,0)=0$, and $a=\Suc b$ gives $n\leq b$ and then
$n\leq\Suc b$ by persistence.

\emph{Antisymmetry.} Fix $a$ and induct on $n$ in
$(a\leq n\land n\leq a)\to a=n$. At $\Suc n$ the equivalence splits
$a\leq\Suc n$ into $a=\Suc n$, or $a\leq n$; in the latter case
$\Suc n\leq a$ gives $n\leq a$, the hypothesis gives $a=n$, and
$\Suc n\leq n$ is a contradiction.

\emph{Transitivity.} Fix $a$, $b$ and induct on $c$ in
$(a\leq b\land b\leq c)\to a\leq c$. At $\Suc c$ the equivalence splits
$b\leq\Suc c$ into $b\leq c$, where the hypothesis and persistence apply, or
$b=\Suc c$, where substitution applies.

For the last statement: if $n\leq a$, the first statement gives
$\operatorname{Cap}(n,a)=\Suc^2n$ directly, or
$\operatorname{Cap}(n,a)=\Suc^2a$, which is $a\leq n$ and gives $a=n$ by
antisymmetry. If $n\not\leq a$, then $a\leq n$ by linearity, which is
$\operatorname{Cap}(n,a)=\Suc^2a$.
\end{proof}

\begin{lemma}[Bounded selection]\label{seed:lem:tf-sel}
For Boolean $c,p,q$ let
$\mu(c,p,q)\defeq(c\land_Bp)\lor_B(\neg_Bc\land_Bq)$, and let
\[
 \Theta(b,c,a,d)\defeq F\bigl(b;\,v.\,
 C(\mu(c,\varepsilon(v,\Suc^2a),\varepsilon(v,\Suc^2d)),v)\bigr),\qquad
 \operatorname{Sel}(b,c,a,d)\defeq w(w(\Theta(b,c,a,d))).
\]
If $a=\Suc^2a_0$ and $a\leq b$, then $c=1\to\operatorname{Sel}(b,c,a,d)=a$.
If $d=\Suc^2d_0$ and $d\leq b$, then $c=0\to\operatorname{Sel}(b,c,a,d)=d$.
\end{lemma}

\begin{proof}
If $c=1$ the body equals $\lambda_a(v)$ for a generic $v$, so
$\Theta(b,c,a,d)=\operatorname{Cap}(b,a)$ by induction on the equation
between the two recursors at a counter variable; similarly
$\Theta(b,c,a,d)=\operatorname{Cap}(b,d)$ if $c=0$. Under the hypotheses
the selected value is $\Suc^2a=\Suc^4a_0$ or $\Suc^2d=\Suc^4d_0$, and two
applications of \eqref{seed:eq:W} remove two successors.
\end{proof}

\subsection{Bounded iteration of finite-case functions}

Write $e(x)\defeq\Suc^2x$. On an input $u=e(x)$ the finite cases
\begin{equation}\label{seed:eq:decode}
 u=2:\ 0;\qquad u=3:\ 1;\qquad u\neq2,3:\ w(w(u))
\end{equation}
return $x$: in the last case $u=\Suc^4x'$ by coverage. For $B=e(b)$ let
$\operatorname{clip}_B(v)\defeq w(w(\operatorname{Cap}(v,B)))$. For
$v\geq2$, \cref{seed:lem:tf-order} and \eqref{seed:eq:W} give
$\operatorname{clip}_B(v)=\min(v,B)$ and
$2\leq\operatorname{clip}_B(v)\leq B$. The term is applied only to
arguments $e(t)$ or to values in $[2,B]$.

\begin{proposition}[Bounded iteration]\label{seed:prop:tf-iter}
Let $g(\vec p)$ and $h(v,\vec p)$ be finite-case functions. For every cap
variable $b$ there is a finite-case function
$\operatorname{It}_b(n,\vec p)$, obtained from terms of the original
language by the cases \eqref{seed:eq:decode}, such that $\mathsf{T}_F$
proves
\[
 \operatorname{It}_b(0,\vec p)=\min(g(\vec p),b),\qquad
 \operatorname{It}_b(\Suc n,\vec p)
 =\min\bigl(h(\operatorname{It}_b(n,\vec p),\vec p),b\bigr),\qquad
 \operatorname{It}_b(n,\vec p)\leq b .
\]
Every occurrence of $F$ in the construction has a body that is a term of the
original language.
\end{proposition}

\begin{proof}
Put $B\defeq e(b)$. Let $g$ have the cases $(A_i,t_i)$. Encode and clip each
candidate as $\operatorname{Cap}(t_i,b)=e(\min(t_i,b))$, which lies in
$[2,B]$ for all inputs, and fold the cases by nested
$\operatorname{Sel}(B,\chi_{A_i},\cdot,\cdot)$ with the default value $2$.
By \cref{seed:lem:tf-sel} this gives a term $a_B(\vec p)$ in $[2,B]$ with
$a_B(\vec p)=e(\min(g(\vec p),b))$.

For the update, compose the cases of $h$ with the cases
\eqref{seed:eq:decode} for its state argument, refine the guards, and encode,
clip and select in the same way. The result is a term $h_B(u,\vec p)$ of
the original language with $2\leq h_B(u,\vec p)\leq B$ for every $u$,
because every candidate is a value of $\operatorname{Cap}(\cdot,b)$, and
\[
 h_B(e(v),\vec p)=e(\min(h(v,\vec p),b)).
\]
All guard tests are Boolean terms; nothing depends on an order of
evaluation.

The seed of the recursor is tied to the body, so the initial value $a_B$ is
entered through two transient states. Let
\[
 H_B(u)\defeq\operatorname{Sel}\Bigl(\Suc^2B,\varepsilon(u,0),B,\;
 \operatorname{Sel}\bigl(\Suc^2B,\varepsilon(u,\Suc B),\Suc^2B,\;
 \operatorname{Sel}(\Suc^2B,\varepsilon(u,\Suc^2B),a_B,h_B(u))\bigr)\Bigr).
\]
All candidates lie in $[2,\Suc^2B]$. The sentinels $\Suc B$ and $\Suc^2B$
are distinct and exceed $B$, so
\[
 H_B(0)=B,\qquad H_B(\Suc B)=\Suc^2B,\qquad H_B(\Suc^2B)=a_B,\qquad
 2\leq u\leq B\to H_B(u)=h_B(u),
\]
and therefore $F(0;u.H_B)=\Suc B$, $F(1;u.H_B)=\Suc^2B$,
$F(2;u.H_B)=a_B$. Put $Y_b(n,\vec p)\defeq F(e(n);u.H_B(u))$.
Quantifier-free induction gives $2\leq Y_b(n,\vec p)\leq B$,
$Y_b(0,\vec p)=a_B$ and
$Y_b(\Suc n,\vec p)=h_B(Y_b(n,\vec p),\vec p)$. Applying the cases
\eqref{seed:eq:decode} to $Y_b$ gives $\operatorname{It}_b$.
\end{proof}

The proposition gives initialized, capped iteration of finite-case
functions whose update sees the state only. Each use below supplies a cap
that is a term of the original language in one variable and a proof that
clipping does not occur on the intended arguments.

\subsection{Arithmetic by finite-case graphs}

All functions introduced in this subsection are finite-case, every induction
is quantifier-free after \cref{seed:lem:tf-cases}, and
\cref{seed:prop:tf-iter} is used whenever an update refers to one of them.

\begin{lemma}[Difference and doubling]\label{seed:lem:tf-diff}
There is a finite-case function $d$ with
\[
 d(x,0)=x,\qquad d(x,\Suc n)=P(d(x,n)),\qquad d(x,n)\leq x,\qquad
 d(\Suc x,\Suc y)=d(x,y),\qquad d(n,n)=0 .
\]
The term $Q(n)\defeq F(n;v.\Suc^2v)$ and the finite-case function
$\operatorname{db}(n)\defeq P(P(P(Q(n))))$ satisfy
\[
 Q(0)=3,\quad Q(\Suc n)=\Suc^2Q(n),\qquad
 \operatorname{db}(0)=0,\quad
 \operatorname{db}(\Suc n)=\Suc^2\operatorname{db}(n),\quad
 Q(n)=\Suc^3\operatorname{db}(n),
\]
and $\operatorname{db}$ is monotone with $n\leq\operatorname{db}(n)$.
\end{lemma}

\begin{proof}
Let $d(x,n)$ be \cref{seed:prop:tf-iter} with initial value $x$, update $P$
and cap $x$. Since $P(v)\leq v$, clipping is never active, which gives the
first three statements. $d(\Suc x,\Suc y)=d(x,y)$ by induction on $y$ with
$x$ fixed: the base is $P(\Suc x)=x$ and the step applies $P$. Then
$d(n,n)=0$ by induction. The equations for $Q$ are \eqref{seed:eq:F}. By
induction $Q(n)\geq3$, so the functional coverage $n=0\lor n=\Suc(P(n))$
cancels three successors and gives the equations for $\operatorname{db}$.
$n\leq\operatorname{db}(n)$ is a quantifier-free induction, and monotonicity
is an induction on the larger argument in
$a\leq b\to\operatorname{db}(a)\leq\operatorname{db}(b)$ with the
equivalence of \cref{seed:lem:tf-order}.
\end{proof}

\begin{lemma}[Addition]\label{seed:lem:tf-add}
There is a finite-case function $+$ with $0+a=a$,
$\Suc n+a=\Suc(n+a)$, commutative, associative and monotone, and with
$n+a<Q(b)$ whenever $n,a\leq b$.
\end{lemma}

\begin{proof}
For a cap $c$ let $A_c(n,a)$ be \cref{seed:prop:tf-iter} with initial value
$a$ and update $\Suc$, and write $s_c(v)\defeq\min(\Suc v,c)$:
\[
 A_c(0,a)=\min(a,c),\qquad A_c(\Suc n,a)=s_c(A_c(n,a)).
\]
For $a\leq c$, induction on the equation gives the seed commutation
$A_c(n,s_c(a))=s_c(A_c(n,a))$: both sides have the same value at $0$ and the
same update. The minimum satisfies $\min(0,c)=0$ and
$\min(\Suc n,c)=s_c(\min(n,c))$ by the order laws, so
$A_c(n,0)=\min(n,c)$ by uniqueness of iteration.

Fix $a\leq c$ and induct on $n$ in $n\leq c\to A_c(n,a)=A_c(a,n)$. The base
is the last identity. If $\Suc n\leq c$, the step applies $s_c$ on the left
and the seed commutation on the right. Since the update increases and $s_c$
preserves order, $A_c$ is monotone in the counter (induction on the larger
counter) and in the initial value within $[0,c]$ (induction on the common
counter).

Fix an input bound $b$ and take the cap $c\defeq Q(b)=\Suc^3
\operatorname{db}(b)$, a term of the original language. Induct on $n$ in
$n\leq b\to A_c(n,n)=\operatorname{db}(n)$: at a successor the seed
commutation gives $A_c(\Suc n,\Suc n)=s_c(s_c(A_c(n,n)))$, and since
$\operatorname{db}(\Suc n)\leq\operatorname{db}(b)<c$ neither step is
clipped, so the value is $\Suc^2\operatorname{db}(n)$. By monotonicity
\[
 n,a\leq b\to A_{Q(b)}(n,a)\leq A_{Q(b)}(b,b)=\operatorname{db}(b)<Q(b),
\]
so there is no clipping on bounded inputs. For the caps of two input bounds
$b$, $b'$, induction on $n$ restricted by $n\leq b,b'$ shows that the
results agree for $a\leq b,b'$: the values at $0$ are $a$ and both steps are
unclipped successors.

Take $b$ to be the first largest of the two arguments, by finitely many
order cases; this defines the finite-case function $n+a$. Coherence across
the caps proves $0+a=a$ and $\Suc n+a=\Suc(n+a)$, bounded commutativity
proves commutativity, and associativity and monotonicity are
quantifier-free inductions with fixed parameters.
\end{proof}

\begin{lemma}[Multiplication and pairing]\label{seed:lem:tf-pair}
Let $E(n)\defeq F(n;v.Q(v))$, so that $E(0)=4$ and $E(\Suc n)=Q(E(n))$.
There are finite-case functions $\cdot$, $\kappa$, $\pi_0$, $\pi_1$ with
\begin{gather*}
 n\cdot a\leq E(Q(b))\ \text{for}\ n,a\leq b,\qquad
 \kappa(x,y)\defeq(x+y)^2+x\leq Q(E(Q(Q(b))))\ \text{for}\ x,y\leq b,\\
 \pi_0(\kappa(x,y))=x,\qquad\pi_1(\kappa(x,y))=y,\qquad
 \pi_0(z)\leq z,\qquad\pi_1(z)\leq z,
\end{gather*}
and the semiring identities for $+$ and $\cdot$.
\end{lemma}

\begin{proof}
$E$ is monotone, $E(n)\geq n$ and
$E(\Suc n)=\operatorname{db}(E(n))+3$. Fix $n,a\leq b$ and apply
\cref{seed:prop:tf-iter} with initial value $0$, update $v\mapsto v+a$ and
cap $E(Q(b))$. For the state $v_j$ at a stage $j\leq b$ the invariant
$v_j\leq E(b+j)$ holds: the base is immediate, and
\[
 v_j+a\leq E(b+j)+b\leq\operatorname{db}(E(b+j))<E(\Suc(b+j)),
\]
using $E(b+j)\geq b$. Since $b+j\leq\operatorname{db}(b)<Q(b)$, the stage
bound is at most the cap, so clipping is inactive. Cap coherence as for
addition gives a finite-case multiplication with the bound $E(Q(b))$; its
equations at $0$ and at a successor give the semiring identities by
quantifier-free induction, in particular $w^2+w<(\Suc w)^2$.

$\kappa$ is finite-case with the stated bound, which is monotone and at
least its argument, and both arguments of a code are at most the code. For
the root, apply \cref{seed:prop:tf-iter} with initial value $0$, the update
that sends $u$ to $\Suc u$ if $(\Suc u)^2\leq z$ and to $u$ otherwise, the
cap $z$ and the counter $z$. The update preserves $u\leq z$ and
$u^2\leq z$, and the quantifier-free invariant ``$u_j=j$ or the update
fixes $u_j$'' shows that the state at $j=z$ is fixed: if it equals $z$,
then $(\Suc z)^2>z$. So the result $r(z)$ satisfies
$r(z)^2\leq z<(\Suc r(z))^2$, and it is unique with this property by
monotonicity of squares. Put
\[
 \pi_0(z)\defeq d(z,r(z)^2),\qquad\pi_1(z)\defeq d(r(z),\pi_0(z)).
\]
The identity $d(x+y,y)=x$ follows by induction on $y$ from
$d(\Suc x,\Suc y)=d(x,y)$ and the addition equations. For
$z=\kappa(x,y)$ and $w\defeq x+y$, the inequalities
$w^2\leq z\leq w^2+w<(\Suc w)^2$ give $r(z)=w$, hence the projection
equations. The projections are total on all arguments and bounded by their
argument, since a difference is bounded by its first argument.
\end{proof}

\begin{theorem}[Primitive recursion]\label{seed:thm:tf-primrec}
For every primitive recursive description $f$ there are a finite-case
function $f^{*}$ in $\mathsf{T}_F$, proofs of the defining equations of $f$
for $f^{*}$, and a term $M_f$ in one variable, monotone and with
$b\leq M_f(b)$, such that $\vec x\leq b\to f^{*}(\vec x)\leq M_f(b)$ is
provable.
\end{theorem}

\begin{proof}
By recursion on $f$. Zero, successor and projections have these data with
the bound $Q$. For a composition $f=h(g_1,\ldots,g_m)$ take a common
refinement of the cases and the bound
$M_h\circ M_{g_1}\circ\cdots\circ M_{g_m}$. Nullary descriptions have
numeral values and constant shifts of $b$ as bounds.

Let $f(0,\vec p)=g(\vec p)$ and
$f(\Suc n,\vec p)=h(n,f(n,\vec p),\vec p)$. Write $K_\kappa$ for the bound
of $\kappa$ in \cref{seed:lem:tf-pair} and put
\[
 \Delta\defeq K_\kappa\circ M_h\circ M_g\circ Q,\qquad
 \Upsilon(t)\defeq F(t;v.\Delta(v)),\qquad
 \beta(b)\defeq\Upsilon(\Suc Q(b)).
\]
Each factor is monotone and at least its argument, and
$\Delta(v)\geq\Suc v$; so $\Upsilon$ is monotone with
$\Upsilon(t)\geq\Suc t$, and $\beta$ is monotone with $\beta(b)\geq b$, by
quantifier-free inductions. The state update is the finite-case function
\[
 J(z,\vec p)\defeq
 \kappa\bigl(\Suc\pi_0(z),\,h^{*}(\pi_0(z),\pi_1(z),\vec p)\bigr),
\]
with the initial value $\kappa(0,g^{*}(\vec p))$. Apply
\cref{seed:prop:tf-iter} to them with the cap $\beta(b)$ and let $z_j$ be
the state at the stage $j$. For $n,\vec p\leq b$ and $j\leq b$,
\begin{equation}\label{seed:eq:prstage}
 z_j\leq\Upsilon(\Suc b+j).
\end{equation}
The seed satisfies
$\kappa(0,g^{*}(\vec p))\leq K_\kappa(M_g(b))\leq\Delta(b)
\leq\Delta(\Upsilon(b))=\Upsilon(\Suc b)$. For the step write
$s\defeq\Upsilon(\Suc b+j)$. Then $s\geq b$, both projections of $z_j$ are
at most $s$, the next counter is at most $\Suc s\leq Q(s)$ and the next
value is at most $M_h(M_g(Q(s)))$, so
$J(z_j,\vec p)\leq\Delta(s)=\Upsilon(\Suc(\Suc b+j))$. Since $j\leq b$, the
argument $\Suc b+j$ is at most $\Suc\operatorname{db}(b)\leq\Suc Q(b)$, so
every stage bound is at most $\beta(b)$ and no clipping occurs. The
invariant \eqref{seed:eq:prstage} is quantifier-free in the extension.

By the projection equations and a quantifier-free induction guarded by
$j\leq b$, the first component of $z_j$ is $j$, and the second component
satisfies the base equation and, for $\Suc j\leq b$, the step equation. Two
caps give the same run up to their common bound on the counter: the seeds
agree and the unclipped updates are identical. Take $b$ to be the first
largest among $n,\vec p$ and define $f^{*}(n,\vec p)\defeq\pi_1(z_n)$. For
any larger input bound, cap coherence identifies this value with the run at
that cap, so $f^{*}(n,\vec p)\leq\beta(b)$ for every valid $b$; put
$M_f\defeq\beta$. Cap coherence proves the defining equations also when a
successor input changes the first largest coordinate.

\Cref{seed:prop:tf-iter} expands every use of $g^{*}$, $h^{*}$, $\kappa$,
$\pi_i$ inside the update into terms of the original language, encoded
cases and bounded selections. There is no evaluator for codes of
descriptions and no recursor over functions defined by graphs.
\end{proof}

\begin{theorem}[Bi-interpretation]\label{seed:thm:tf-biint}
$\mathsf{T}_F$ and $\iPRA$ are bi-interpretable, with whole domains, actual
equality, the logic translated identically and the identity as isomorphism
of both composites. They are equiconsistent.
\end{theorem}

\begin{proof}
In $\iPRA$ keep $0$, $\Suc$ and equality, interpret $\varepsilon$ as the
primitive recursive characteristic of equality, and each recursor with body
$h$ by primitive recursion with initial value $\Suc(h(0))$ and a step $h$
that ignores the counter. All axioms of $\mathsf{T}_F$ and its induction
rule hold. Conversely interpret each description by the finite-case
function of \cref{seed:thm:tf-primrec}; totality, single-valuedness, the
defining equations and the quantifier-free induction instances are proved in
$\mathsf{T}_F$ by \cref{seed:lem:tf-cases}, and injectivity and separation
are axioms.

On the side of $\iPRA$, recursion on a description reduces the
identification of a function with its double interpretation to uniqueness
from the base and step equations, which is a quantifier-free induction. On
the side of $\mathsf{T}_F$, $0$ and $\Suc$ are unchanged; the two
characteristics of equality are two-valued and reflect equality, so they
agree by cases; and for a recursor with body $h$, the earlier stage
identifies the body, both functions satisfy \eqref{seed:eq:F} for it, and
induction on the equation between them identifies them, with
\cref{seed:lem:tf-cases} where graph-defined functions occur. Induction on
formulas extends this to all formulas, and falsity is preserved.
\end{proof}

%======================================================================
\section{Three observations on the native core}
\label{seed:sec:core}

The packages of this chapter contain three results about the calculus
$\Nat$ on iterator terms that do not concern a seeded recursor. They are
recorded here from \pkg{coverage\_and\_kernel}, \pkg{transient\_seed} and
\pkg{monus\_and\_neutral\_major}. The second coincides with
\cref{kern:prop:compoundonly}(ii) and is recorded by a reference; the first
and the third differ from the results of \cref{kern:sec:onesided} and
\cref{claus:sec:prefix} with similar names, and the remarks below relate
them.

\subsection{One native binder over a literal equation}

In this subsection formulas are read in $\iPRA$ through
\cref{base:sec:reading,base:thm:bridge}, and $\sel$, the predecessor $P$ of
\cref{base:thm:pred}, the characteristic terms of \cref{base:thm:reflect}
and the pairing of \cref{base:thm:primrec} are available. A formula is
called $\Sigma_2^{\mathrm{qf}}$ if it has the form
$\exists x\,\forall y\,\theta(x,y,\vec p)$ with $\theta$ quantifier-free,
and $\Pi_1^{\mathrm{qf}}$ if it has the form $\forall y\,\theta(y,\vec p)$.
These are classes of prenex formulas; no prenex normal form theorem for
intuitionistic formulas is used.

\begin{proposition}[One binder]\label{seed:prop:onebinder}
Up to provable equivalence, the native formulas
$\D{x}{z}E(x,z,\vec p)$ with $E$ a literal equation are exactly the
$\Sigma_2^{\mathrm{qf}}$ formulas, and the native formulas
$\Q{z}E(z,\vec p)$ with $E$ a literal equation are exactly the
$\Pi_1^{\mathrm{qf}}$ formulas.
\end{proposition}

\begin{proof}
The reading of $\D{x}{z}E$ is
$\exists x\,(E(x,0)\to\forall z\,E(x,z))$, which is equivalent to
$\exists x\,\forall z\,(E(x,0)\to E(x,z))$ because $z$ does not occur in
$E(x,0)$; the matrix is quantifier-free. Without the existential quantifier
this is the statement for $\Q{z}$.

Conversely let $r(x,y,\vec p)$ be a characteristic term with
$r=0\leftrightarrow\theta$, and let $E(x,z,\vec p)$ be the literal equation
\[
 \sel\bigl(z,0,r(x,P(z),\vec p)\bigr)=0 .
\]
Its instance at $z=0$ is provable, and its instance at $z=\Suc y$ is
equivalent to $r(x,y,\vec p)=0$ by $P(\Suc y)=y$. By coverage,
$\forall z\,E(x,z,\vec p)\leftrightarrow\forall y\,\theta(x,y,\vec p)$, and
the provable antecedent $E(x,0,\vec p)$ drops out of the reading of
$\D{x}{z}E$.
\end{proof}

\begin{corollary}\label{seed:cor:onebinder}
If rule \rulename{QI} is extended to matrices that consist of one
$\D{x}{y}$ over a literal equation, the extended calculus derives every
induction instance for $\Sigma_2^{\mathrm{qf}}$ formulas. If it is extended
to matrices $\Q{y}E$ with $E$ a literal equation, it derives every induction
instance for $\Pi_1^{\mathrm{qf}}$ formulas.
\end{corollary}

\begin{proof}
Let $B(n,\vec p)$ be the native formula of \cref{seed:prop:onebinder} for
the given formula, and let $\Gamma\proves B(0,\vec p)$ and
$\Gamma,B(z,\vec p)\proves B(\Suc z,\vec p)$. The extended rule gives
$\Q{n}B(n,\vec p)$, and \rulename{QE} with the base gives $B(t,\vec p)$ for
every term $t$. The equivalence of the proposition transports the
derivation to the given formula.
\end{proof}

A guard ``at most one binder'' is therefore not equivalent to the guard
``literal equation'' unless an elimination or conservativity theorem for the
additional induction instances is proved. No statement on the strength of
the extended calculus relative to $\iPRA$ is made.

For a native formula let $\delta(s=t)\defeq0$ and
$\delta(\Q{y}A)\defeq\delta(\D{x}{y}A)\defeq1+\delta(A)$. For a fresh $z$ let
$G_zA$ be $A$ with every equation $s=t$ replaced by
$\sel(z,0,s)=\sel(z,0,t)$. Then $\delta(G_zA)=\delta(A)$, $G_0A$ is
provable, and $G_{\Suc u}A\leftrightarrow A$: for equations by computation,
and for a binder by induction on the formula, with the witness $0$ for
$\D{x}{y}$.

\begin{theorem}[Alternating prefixes]\label{seed:thm:prefix}
For every $k\geq0$, every formula
$\exists x_1\forall y_1\cdots\exists x_k\forall y_k\,
\theta(\vec x,\vec y,\vec p)$ with $\theta$ quantifier-free is provably
equivalent in $\iPRA$ to the reading of a native formula of depth
$\delta\leq k$. Consequently an extension of \rulename{QI} to all matrices
of depth at most $k$ derives induction for this class.
\end{theorem}

\begin{proof}
For $k=0$ take the characteristic equation. If $A(x,y,\vec p)$ translates
the remaining prefix at depth at most $k-1$, put
$\mathcal B(A)\defeq\D{x}{z}\,G_zA(x,P(z),\vec p)$. Its matrix at $z=0$ is
provable and its matrix at $z=\Suc y$ is equivalent to $A(x,y,\vec p)$, so
by coverage $\mathcal B(A)\leftrightarrow\exists x\,\forall y\,A(x,y,\vec p)$,
and $\delta(\mathcal B(A))\leq k$. The induction statement is the argument
of \cref{seed:cor:onebinder}.
\end{proof}

Blocks of several variables are packed by pairing, and dummy blocks can be
added because the domain is inhabited; so the inclusion covers the prenex
$\Sigma_{2k}$ formulas and, for $k\geq1$, the prenex $\Pi_{2k-1}$ formulas.
At depth greater than one only this inclusion is claimed: the reading of a
nested native formula contains implications with quantified antecedents, and
moving quantifiers through them is not valid intuitionistically.

\begin{remark}[Relation to the prefix representation]
\label{seed:rem:prefix-claus}
\Cref{claus:lem:prefix,claus:thm:prefix} concern a different statement with
a similar name: every $Q$-free native formula is, uniquely and stably under
substitution, a record $D^k(s=t)$, and a clause calculus on such records is
equivalent to $\Nat$. That is a statement about syntax and derivations.
\Cref{seed:prop:onebinder,seed:thm:prefix} concern what native formulas
express under the reading in $\iPRA$. The two fit together: the formulas
built in the proof of \cref{seed:thm:prefix}, and those for the
$\Sigma_2^{\mathrm{qf}}$ case of \cref{seed:prop:onebinder}, use $\D{}{}$, an
equation and the gate $G_z$ only, so they are $Q$-free and are records $D^{j}(s=t)$ with $j\leq k$ in the sense of
\cref{claus:lem:prefix}.
\end{remark}

\subsection{The case rule at a literal equation cannot be deleted}

\begin{remark}[The atomic case introduction]\label{seed:prop:atomic-case}
\pkg{transient\_seed} considers the calculus on iterator terms with paired
introduction by cases and the progress rule on $\D{x}{y}$ of
\cref{seed:rem:succ-rules}, with the introduction restricted so that its
local assumption is $E(t,z)$ whenever the matrix is a literal equation $E$,
and $B(t,0)$ otherwise, and shows that the restricted calculus does not
derive $\D{x}{y}(y=\Suc 0)$, which the unrestricted calculus derives. At a
literal equation the restricted introduction is the progress rule
\rulename{DP}, so the restricted calculus is that of
\cref{kern:prop:compoundonly}(ii) with transport in a formula template, and
the statement and its proof, by the reading of an atomic $\D{x}{y}E$ as
$\exists x\,\forall y\,(E(x,y)\to E(x,\Suc y))$, are given there.
\Cref{guard:prop:strict} is the same statement for the root-determined
introduction $\rulename{DI}_{\mathrm{root}}$, and
\cref{kern:prop:nodi,kern:prop:compoundonly}(i) are the companion results
for dropping \rulename{DI} altogether and for keeping \rulename{QI} in place
of progress.

Under the standard reading the same formula is true, since its zero instance
is false. Hence the two readings of a $\D{x}{y}$ over an equation, closure
and progress, cannot be exchanged by padding the equation into a compound
formula without losing one direction of the formula round trip.
\end{remark}

\subsection{Transport from the literal zero alone}

\begin{proposition}\label{seed:prop:zero-source}
In the calculus on iterator terms with formula transport
(\cref{seed:rem:succ-rules}) and no reflexivity axiom, replace transport by
the rule: from $\Gamma\proves0=t$ and $\Gamma\proves F(0)$ infer
$\Gamma\proves F(t)$. The resulting calculus does not derive $0=0$.
\end{proposition}

\begin{proof}
Interpret terms in $\{0,1\}$ with $\Suc 0=\Suc 1=1$ and every iterator
term equal to $1$, read an equation $s=t$ as ``the value of $s$ is $1$'',
and read $\Q{y}$, $\D{x}{y}$ by the clauses of \cref{base:sec:reading}.
Every computation equation holds, because its left side is an iterator. No
equation $0=t$ holds, so the restricted rule is sound. Every element is $0$
or a successor, and induction holds for every predicate on the domain: the
base gives the value at $0$ and the step at $0$ gives the value at $1$. So
\rulename{DI} by cases and the progress rule are sound, as are
\rulename{Hyp}, \rulename{QE} and \rulename{DE}. The equation $0=0$ does not
hold.
\end{proof}

The relation used in the proof is not equality. With unrestricted transport
$0=0$ follows from the computation $\It{0}{0}{v}{0}=0$ used as both premises
with the template $w=0$. The proposition does not apply to a calculus that
has $0=0$ as an axiom or as an initial context entry.

\begin{remark}[The opposite restriction of the source]
\label{seed:rem:zero-source}
The proposition restricts transport to the source $0$, which is the
complement of the restriction $s\not\equiv0$ of \cref{kern:sec:onesided}.
That restriction is harmless: reflexivity and the unrestricted rule
\rulename{Tr} are derivable from \rulename{I0} and $\rulename{Tr}^{-}$
(\cref{kern:lem:oseq}), and for transport in a formula template the calculus
$\mathsf Z$, in which a source $0$ is served by \rulename{QE}, is equivalent
to $\Nat$ (\cref{kern:thm:dispatch}). The restriction $s\not\equiv0$ is
harmful only inside the strictly guarded composition of
\cref{guard:cor:obstruction}, where a zero-left hypothesis is inert. Thus the
sources other than $0$ suffice for equality in $\Nat$, and the source $0$
alone does not.
\end{remark}

%======================================================================
\section{Checker sizes}
\label{seed:sec:sizes}

All eight packages report sizes of complete standalone checkers under the
protocol of \cref{base:sec:measure}. \Cref{seed:tab:sizes} lists them as
stated, with the control each package measures against. Only the programs
of the first three packages that are marked $R$ check the successor-seeded
calculus $\mathsf{N}_R$. All other programs check the calculus $\Nat$ on
iterator terms; no package reports a checker for
$\mathsf{T}_{B\varepsilon}$, $\mathsf{T}_{Bd}$, $\mathsf{N}_{Bd}$,
$\mathsf{T}_F$ or $\mathsf{T}_F^{-}$, and the sizes below are not sizes of
those theories.

\begin{longtable}{@{}llrrr@{}}
\caption{Checker sizes as stated by the packages (tokens, AST nodes,
normalized bytes).}\label{seed:tab:sizes}\\
\toprule
Package & Program & Tokens & AST & Bytes\\
\midrule
\endfirsthead
\toprule
Package & Program & Tokens & AST & Bytes\\
\midrule
\endhead
\bottomrule
\endfoot
\pkg{scalar\_successor} & \texttt{checker.py} ($R$) & 279 & 319 & 686\\
 & \texttt{checker\_bytes.py} ($R$) & 279 & 321 & 681\\
 & \texttt{checker\_ast.py} ($R$) & 282 & 314 & 708\\
 & control: generated assumptions, main & 281 & 318 & 685\\
 & control: generated assumptions, AST profile & 283 & 315 & 710\\
 & control: guarded elimination, token profile & 281 & 325 & 739\\
\midrule
\pkg{transient\_seed} & \texttt{checker\_sr\_bytes.py} ($R$) & 284 & 322 & 681\\
 & \texttt{checker\_bytes.py} (iterator) & 282 & 320 & 679\\
 & \texttt{checker\_structure.py} (iterator) & 280 & 311 & 705\\
 & \texttt{checker\_joint.py} (iterator, predecessor packets) & 280 & 314 & 684\\
 & control: generated assumptions, main & 281 & 318 & 685\\
 & control: generated assumptions, AST profile & 283 & 315 & 710\\
\midrule
\pkg{successor\_seed\_completion} & \texttt{checker\_direct.py} ($R$) & 283 & 320 & 687\\
 & \texttt{checker.py} ($R$) & 284 & 322 & 681\\
 & \texttt{checker\_ast.py} ($R$) & 282 & 313 & 707\\
 & \texttt{checker\_ast\_total.py} ($R$) & 287 & 314 & 711\\
 & control: prior generated-assumption checker & 281 & 318 & 685\\
 & control: prior AST-oriented checker & 283 & 315 & 710\\
\midrule
\pkg{bounded\_initializer} & \texttt{checker.py} & 266 & 310 & 612\\
 & \texttt{checker\_compact.py} & 266 & 312 & 604\\
 & \texttt{checker\_ast.py} & 267 & 298 & 639\\
 & control: byte leader & 269 & 311 & 616\\
 & control: AST and token leader & 268 & 300 & 649\\
\midrule
\pkg{monus\_and\_neutral\_major} & \texttt{checker.py} & 263 & 307 & 550\\
 & \texttt{checker\_bytes.py} & 263 & 309 & 542\\
 & \texttt{checker\_joint.py} & 261 & 295 & 578\\
 & control: certified zero, explicit list test & 264 & 308 & 566\\
 & control: certified zero, compact list test & 264 & 310 & 558\\
 & control: certified zero, generic traversal & 262 & 296 & 598\\
\midrule
\pkg{selfseed\_scalar} & \texttt{local\_argument.py} & 267 & 309 & 562\\
 & \texttt{local\_argument\_bytes.py} & 267 & 311 & 554\\
 & \texttt{local\_loop\_joint.py} & 263 & 291 & 642\\
 & \texttt{local\_loop\_joint\_separate.py} & 262 & 293 & 641\\
 & control: \texttt{baseline\_checker.py} & 264 & 308 & 566\\
 & control: \texttt{baseline\_checker\_bytes.py} & 264 & 310 & 558\\
 & control: \texttt{baseline\_checker\_joint.py} & 262 & 296 & 598\\
\midrule
\pkg{scope\_and\_coverage} & \texttt{checker\_compact.py} & 266 & 310 & 529\\
 & \texttt{checker\_explicit.py} & 266 & 308 & 537\\
 & control: neutral major, global floor (byte incumbent) & 263 & 309 & 542\\
 & control: guarded major, lexical floor & 267 & 311 & 545\\
 & control: neutral major, lexical floor, conditional dispatch & 266 & 310 & 535\\
\midrule
\pkg{coverage\_and\_kernel} & \texttt{checker.py} & 266 & 308 & 539\\
 & \texttt{checker\_bytes.py} & 266 & 310 & 531\\
 & \texttt{checker\_ast.py} & 263 & 289 & 613\\
 & control: byte leader & 263 & 309 & 542\\
 & control: token leader & 261 & 295 & 578\\
 & control: AST leader & 263 & 291 & 642\\
\end{longtable}

In every package the minima over the three coordinates are attained by
different programs, and each package states that no single one of its
programs is below all three coordinates of its control, with one exception:
\texttt{checker\_joint.py} of \pkg{transient\_seed}, at 280/314/684, is
below 281/315/685, the coordinatewise minimum of its two controls. The
packages state that these minima are sizes attained in a fixed corpus of
earlier checkers and not lower bounds.

\paragraph{Successor-seeded checkers.}
The $R$ programs derive from the checker with generated assumptions of an
earlier draft: a certificate is a JSON array whose last element is the
closed goal and whose earlier four-field packets are consumed from the end;
a computation equation enters as a local assumption
(\cref{seed:rem:succ-rules}). An arithmetic packet supplies the data
$[n,a,h]$ and a flag, and the checker itself builds the counter $0$ or
$\Suc n$ of the generated equation, so no shape test on a supplied counter
is needed. The same change applied to the iterator gives the matched
iterator programs of \pkg{transient\_seed}, whose joint profile uses a
changed packet format for predecessor computations. The AST profiles accept
additional raw aliases and reject malformed nonempty strings only through
the recursion limit of the interpreter; \texttt{checker\_ast\_total.py}
rejects them directly. \pkg{transient\_seed} ships kernel certificates for
the Boolean decoder (45 commands) and for a restricted implication (149
commands); \pkg{successor\_seed\_completion} states that the full
predecessor derivation has not been emitted as a kernel certificate.

\paragraph{Iterator checkers in the other packages.}
The programs of \pkg{bounded\_initializer} compose an arity encoding of the
raw syntax with a predecessor supply. Those of
\pkg{monus\_and\_neutral\_major} implement its neutral-major certificate
protocol for $\Nat$, in which every instruction proves a major premise and a
derivable truth serves where none is needed; the protocol is not described
in this chapter. Those of
\pkg{selfseed\_scalar}, \pkg{scope\_and\_coverage} and
\pkg{coverage\_and\_kernel} compose lexical proof scopes, the arity
encoding, the certified-zero dispatch and, in the last two, neutral majors;
the loop variants keep contexts as immutable tuples. \pkg{selfseed\_scalar}
records two failed variants: a loop with in-place mutation of the context
list accepts a certificate of $0=\Suc 0$ by leaking an assumption of a
subproof, and a variant with opaque objects as fresh parameters accepts a
certificate of $0=\Suc 0$ because a floating-point leaf $3.0$ is treated as
a rigid parameter while comparing equal to the bound index $3$; the stricter
repair of the latter loses the canonical certificate of universal
reflexivity.

%======================================================================
\section{Further questions}
\label{seed:sec:questions}

No result of this chapter has a Lean formalization recorded in the
repository; the packages present written proofs with finite numerical
checks of the defining terms on standard numbers.

\begin{question}[Injectivity for the self-seeded recursor]
Does $\mathsf{T}_F^{-}$ prove $\Suc x=\Suc y\to x=y$? By
\cref{seed:lem:tf-retract} this is equivalent to the closure under successor
of the fixed points of $n\mapsto P(\Suc n)$, and by
\cref{seed:lem:tf-orbit} to the derivation of $p_*(\Suc^2m)=\Suc m$ from
$\Suc(p_*(\Suc^2m))=\Suc^2m$. \pkg{selfseed\_scalar},
\pkg{scope\_and\_coverage} and \pkg{coverage\_and\_kernel} leave it open. A
countermodel has to satisfy all staged recursor equations and all instances
of induction for literal equations, together with a successor collision; a
branching successor graph without cycles is not enough. Neither the
bi-interpretability of $\mathsf{T}_F^{-}$ with $\iPRA$ nor its failure is
established.
\end{question}

\begin{question}[Uniform rigidity]
\Cref{seed:prop:tf-rigid,seed:thm:tf-acyclic} are families of derivations
indexed by a fixed integer. Is there a single derivation with a variable
index, without a general compiler for primitive recursion in
$\mathsf{T}_F^{-}$?
\end{question}

\begin{question}[Shorter predecessor terms]
Is there a substantially shorter predecessor term for the successor-seeded
recursor than those of \cref{seed:thm:succ-pred,seed:rem:succ-three}, for
instance a reset cycle that needs no Boolean comparator? The three packages
prove existence and correctness on the whole domain and make no claim on
size.
\end{question}

\begin{question}[Other seed maps]
For which maps $\sigma$ in place of the successor does the recursor with
$R(0,a;u.h)=\sigma(a)$ recover arbitrary initialization from its two
equations and induction for literal equations? A criterion has to provide a
proved inverse and not assume one. Related: is there a general class of
finite reset mechanisms for which the transient-seed method works, and is
the correction of a capped predecessor at finitely many arguments an
instance of a general repair principle?
\end{question}

\begin{question}[Cancellation on finite-valued terms]
In $\mathsf{N}_R$ and $\mathsf{N}_{Bd}$ a restricted implication and
contextual case analysis are obtained from cancellation of the successor on
two-valued terms, before a selector exists
(\cref{seed:lem:succ-cases,seed:lem:nbd-cases}). Which other weak term
languages admit this route to the full logical interface?
\end{question}

\begin{question}[Tied seeds without an additional operation]
\Cref{seed:thm:tied-bound} excludes terms for addition, selection and
pairing. It does not exclude interpretations of $\iPRA$ by relations,
quotients or several coordinates in a theory on the terms
$x\mid0\mid B(n;u,v.h)$ alone; \cref{seed:thm:tied-biint} needs
$\varepsilon$ and \cref{seed:cor:tbd-biint} needs $d$ and the coverage
axiom. An invariant that is stable under such interpretations is not known
from the packages of this chapter.
\end{question}

\begin{question}[Fewer axioms for the equality test]
$\mathsf{T}_{Bd}$ replaces the range and reflection axioms of
$\mathsf{T}_{B\varepsilon}$ by two recursive equations, at the price of
coverage as an axiom. Can range and reflection be replaced by recursive
equations from which coverage is still derived? \pkg{monus\_and\_neutral\_major}
investigates a basis with three equations for the minimum, in which
$\operatorname{parity}(\min(2x,2y+1))$ has the standard meaning of a
comparator; the derivation of the equations of $d$ and of reflection on the
whole domain from that basis, without induction over a shifted parameter, is
left open there.
\end{question}

\begin{question}[Internal majorants and sizes of the compilers]
The converse direction of \cref{seed:thm:tbe-exact} produces a single term
for a function with a majorant that is valid in the standard model; a proof
in $\mathsf{T}_{B\varepsilon}$ that this term equals the graph-defined
function needs a proof of the majorant in the theory. Bounds on the size of
the terms and proofs produced by
\cref{seed:thm:tbe-realize,seed:lem:tbe-graph,seed:prop:tf-iter,%
seed:thm:tf-primrec}, on the branching of the case expansions for nested
terms, and on the cost of translating the graph-based proofs into a native
calculus are not given.
\end{question}

\begin{question}[Extended progress]
\Cref{seed:cor:onebinder,seed:thm:prefix} show that an extension of
\rulename{QI} to matrices with binders derives induction for prenex
classes. Whether such an extension is conservative over $\Nat$, or strictly
stronger, is not decided; at depth greater than one only the inclusion of
\cref{seed:thm:prefix} is proved.
\end{question}

\begin{question}[Variants of the two obstructions]
\Cref{kern:prop:compoundonly}(ii), in the form of
\cref{seed:prop:atomic-case}, refutes the deletion of case introduction at
literal equations for the two displayed readings; it does not exclude a
calculus in which a $\D{x}{y}$ over an equation is an auxiliary judgment and
the object formulas are padded. \Cref{seed:prop:zero-source} does not
exclude transport from the literal zero together with a reflexivity axiom.
Neither variant is decided in the packages of this chapter; calculi with a
single introduction rule are the subject of \pkg{witness\_rule} and
\pkg{guarded\_core} (\cref{kern:sec:witness,guard:prop:strict}), and
restrictions on the source of transport of \pkg{source\_dispatch}
(\cref{kern:thm:dispatch}).
\end{question}

\begin{question}[Checkers for the new bases]
No package measures a complete checker for $\mathsf{T}_{B\varepsilon}$,
$\mathsf{T}_{Bd}$, $\mathsf{N}_{Bd}$ or $\mathsf{T}_F$, in which the
equality test or the subtraction equations, the ordinary logic where it is
used, and the validation of the staged bodies are all charged. Whether any
of these bases, or $\mathsf{N}_R$, admits a single checker below the
controls of \cref{seed:tab:sizes} in all three coordinates is open, as are
lower bounds for a fixed checker language and input grammar.
\end{question}

\begin{question}[Formalization]
The packages name as first targets for a proof assistant: the contextual
case lemmas (\cref{seed:lem:succ-cases,seed:lem:nbd-cases}), the factor map
identities and the predecessor (\cref{seed:lem:succ-factor,%
seed:thm:succ-pred}), the gated complement (\cref{seed:lem:tbd-gc}), cap
coherence and finite-case elimination
(\cref{seed:thm:tbe-realize,seed:lem:tbe-graph,seed:lem:tf-cases}), and the
coverage and rigidity proofs of \cref{seed:thm:tf-coverage,%
seed:prop:tf-rigid}.
\end{question}

